\RequirePackage{fix-cm}
\documentclass[11pt,a4paper,oneside]{report}
\usepackage[T1]{fontenc}
\usepackage[utf8]{inputenc}
\usepackage{amsmath,amsthm,mathtools}
\usepackage{newtxtext,newtxmath}
\usepackage[margin=27mm,headheight=15pt,headsep=7mm,footskip=12mm]{geometry}
\usepackage{microtype}
\usepackage{booktabs,longtable,array,tabularx}
\usepackage{xcolor,listings,enumitem,needspace}
\usepackage{fancyhdr,titlesec,etoolbox,aliascnt}
\usepackage[unicode,pdfencoding=auto,psdextra]{hyperref}
\usepackage[nameinlink,noabbrev]{cleveref}
\definecolor{ink}{HTML}{162D43}
\definecolor{muted}{HTML}{53616D}
\definecolor{papergray}{HTML}{F5F6F7}
\hypersetup{colorlinks=true,linkcolor=ink,citecolor=ink,urlcolor=ink,
 pdftitle={Quantified CUE: A Set-Theoretic Constraint Calculus},
 pdfauthor={Aram Hăvărneanu},
 pdfsubject={Working technical proposal: quantification, function types, dependency, and abstraction},
 pdfkeywords={CUE, semantic subtyping, hyperdoctrine, polymorphism, existential types}}
\urlstyle{same}
\setlength{\parindent}{0pt}
\setlength{\parskip}{3pt plus 1pt minus 1pt}
\AtBeginDocument{
 \setlength{\abovedisplayskip}{7pt plus 2pt minus 2pt}
 \setlength{\belowdisplayskip}{7pt plus 2pt minus 2pt}
 \setlength{\abovedisplayshortskip}{3pt plus 1pt}
 \setlength{\belowdisplayshortskip}{5pt plus 1pt minus 1pt}
}
\setlength{\emergencystretch}{2em}
\setcounter{secnumdepth}{2}
\setcounter{tocdepth}{1}
\setlist{nosep,leftmargin=1.6em}
\pagestyle{fancy}
\fancyhf{}
\fancyhead[L]{\small\textsc{Quantified CUE}}
\fancyhead[R]{\small\nouppercase{\leftmark}}
\fancyfoot[C]{\thepage}
\renewcommand{\headrulewidth}{0.3pt}
\renewcommand{\chaptermark}[1]{\markboth{\thechapter\quad #1}{}}
\titleformat{\chapter}[display]{\normalfont\Large\bfseries\color{ink}}{\chaptertitlename\ \thechapter}{5pt}{\LARGE}
\titlespacing*{\chapter}{0pt}{0pt}{17pt}
\titleformat{\section}{\large\bfseries\color{ink}}{\thesection}{0.6em}{}
\titleformat{\subsection}{\normalsize\bfseries}{\thesubsection}{0.6em}{}
\titlespacing*{\section}{0pt}{12pt}{5pt}
\titlespacing*{\subsection}{0pt}{9pt}{3pt}
\pretocmd{\section}{\Needspace{9\baselineskip}}{}{}
\pretocmd{\subsection}{\Needspace{7\baselineskip}}{}{}
\lstdefinelanguage{CUE}{
 morekeywords={forall,exists,in,func,let,for,if,package,import,seal,with,open,as,extern,match,case,return},
 sensitive=true,morecomment=[l]{//},morecomment=[s]{/*}{*/},morestring=[b]",
 alsoletter={\#}}
\lstset{language=CUE,basicstyle=\ttfamily\fontsize{8.5}{10.4}\selectfont,
 keywordstyle=\bfseries\color{ink},commentstyle=\color{muted},stringstyle=\color{ink},
 columns=fullflexible,keepspaces=true,showstringspaces=false,tabsize=4,
 breaklines=true,breakatwhitespace=false,upquote=true,
 backgroundcolor=\color{papergray},frame=single,rulecolor=\color{black!13},
 framerule=0.3pt,framesep=5pt,xleftmargin=5pt,xrightmargin=5pt,
 aboveskip=6pt,belowskip=7pt}
\newcommand{\code}[1]{\texttt{\detokenize{#1}}}
\newcommand{\sem}[1]{\llbracket #1\rrbracket}
\newcommand{\Pow}{\mathcal P}
\newcommand{\Env}{\operatorname{Env}}
\newcommand{\FV}{\operatorname{FV}}
\newcommand{\App}{\operatorname{App}}
\newcommand{\Run}{\operatorname{run}}
\newcommand{\Ok}{\mathsf{Ok}}
\newcommand{\Reject}{\mathsf{Reject}}
\newcommand{\Raise}{\mathsf{Raise}}
\newcommand{\Type}{\mathsf{Type}}
\newcommand{\Int}{\mathbb Z}
\newcommand{\Num}{\mathsf{Number}}
\newcommand{\Str}{\mathsf{String}}
\newcommand{\Bool}{\mathsf{Bool}}
\newcommand{\List}{\mathsf{List}}
\newcommand{\Nat}{\mathbb N}
\newcommand{\U}{\mathcal U}
\DeclareMathAlphabet{\profilecal}{OMS}{cmsy}{m}{n}
\newcommand{\Uone}{\ensuremath{\profilecal{U}_{1}}}
\newcommand{\Sone}{\ensuremath{\profilecal{S}_{1}}}
\newcommand{\Shigher}{\ensuremath{\profilecal{S}_{H}}}
\newcommand{\Dep}{\ensuremath{\profilecal{D}}}
\newcommand{\Abs}{\ensuremath{\profilecal{A}}}
\newcommand{\Static}{\ensuremath{\profilecal{K}}}
\newcommand{\Comp}{\mathsf{Comp}}
\newcommand{\Fail}{\mathsf{fail}}
\newcommand{\Pending}{\mathsf{pending}}
\pdfstringdefDisableCommands{%
 \def\Uone{U1}%
 \def\Sone{S1}%
 \def\Shigher{SH}%
 \def\Dep{D}%
 \def\Abs{A}%
 \def\Static{K}%
}
\newcommand{\Rel}{\mathsf{Rel}}
\newcommand{\todoob}{\textbf{Implementation obligation.}\ }
\newtheorem{definition}{Definition}[chapter]
\newaliascnt{proposition}{definition}
\newtheorem{proposition}[proposition]{Proposition}
\aliascntresetthe{proposition}
\newaliascnt{theorem}{definition}
\newtheorem{theorem}[theorem]{Theorem}
\aliascntresetthe{theorem}
\newaliascnt{lemma}{definition}
\newtheorem{lemma}[lemma]{Lemma}
\aliascntresetthe{lemma}
\theoremstyle{remark}
\newaliascnt{remark}{definition}
\newtheorem{remark}[remark]{Remark}
\aliascntresetthe{remark}
\newcounter{example}[chapter]
\renewcommand{\theexample}{\thechapter.\arabic{example}}
\newcommand{\example}[2]{\refstepcounter{example}\Needspace{5\baselineskip}\par\textbf{Example \theexample\ --- #1.}\ \textit{#2}\par\nopagebreak[4]}
\newcommand{\decision}[1]{\par\textbf{Design rule.} #1\par}
\newcolumntype{Y}{>{\raggedright\arraybackslash}X}
\newcolumntype{P}[1]{>{\raggedright\arraybackslash}p{#1}}
\begin{document}
\hypersetup{pageanchor=false}
\begin{titlepage}
\vspace*{16mm}
{\small\textsc{Working technical proposal}\par}
\vspace{12mm}
{\fontsize{31}{35}\selectfont\bfseries\color{ink}Quantified CUE\par}
\vspace{5mm}
{\Large A Set-Theoretic Constraint Calculus\par}
\vspace{3mm}
{\large Polymorphic functions, dependent specifications,\par
and representation-independent modules\par}
\vspace{17mm}
{\large Aram H\u{a}v\u{a}rneanu\par}
\vspace{2mm}
{27 September 2026\par}
\vfill
\begin{minipage}{\textwidth}
\textbf{Abstract.}
CUE descriptions are interpreted as predicates on joint environments of possible
values. Existential projection explains instance formation and hidden witnesses;
universal projection adds obligations over arbitrary values and semantic types.
Both are adjoints to substitution, and both preserve refinement. Function types
are universal predicates on an operational application relation. Their variance,
intersection laws, and polymorphic instances follow from that interpretation.
The proposal defines quantified declarations, legal scope extrusion, a
conservative data-language embedding, and three language profiles. A worked
catalogue covers calling conventions, higher-order programming, dependent
invariants, and abstract modules. Existential modules use sealing and
equality-respecting simulations to support representation-independent evolution.
The reference model establishes refinement, function soundness, and
representation-independent replacement under explicit hypotheses. A mandatory, resource-bounded static kernel checks an annotated polymorphic
structural fragment. Every annotation is discharged by a static derivation. Explicit source
assertions retain their runtime checks and contribute refinements to those
derivations; annotation checking leaves executable code unchanged. Relevant
consistency rejects explicit interfaces whose declared input regions have
statically refuted successful results, independently of denotational inhabitation.
Function arrows express partial correctness;
constraint failure is permitted without making static typing optional.
Impredicative quantification uses a single set-sized family of predicates.
\end{minipage}\par
\vspace{8mm}
{\small Comparison baseline: CUE proposal 4484 and its theory addendum,\par
repository revision \texttt{a949162aedb8}, inspected 21 September 2026.\par}
\end{titlepage}
\pagenumbering{roman}
\hypersetup{pageanchor=true}
\tableofcontents
\clearpage
\chapter*{Conventions}
\addcontentsline{toc}{chapter}{Conventions}
Equations and proofs use the semantic reference calculus. Code listings use
the proposed surface language and carry profile or extension labels. An expected
contradiction is written \code{_|_}; an unresolved proof or evaluation obligation
is written \emph{incomplete}. A blocked static judgment is a compiler
diagnostic, not the semantic value \code{_|_}. Relevant consistency is the
static admissibility condition for explicit interface descriptions. Its failure
blocks an interface even when the underlying function predicate is inhabited.

\begin{tabularx}{\textwidth}{@{}lY@{}}
\toprule
Notation & Meaning\\\midrule
$V$ & A fixed set of operational inhabitants, including data and closures.\\
$\rho\in\Env(\Gamma)$ & An assignment to the free variables of context $\Gamma$.\\
$\sem{T}_\rho\subseteq V$ & The inhabitants of description $T$ at $\rho$.\\
$P\le Q$ & Inclusion of predicates, hence refinement toward more information.\\
$\top,\bot$ & The full subject set and the empty predicate.\\
$\U=\Pow(V)$ & The set of semantic type interpretations.\\
\Uone{} & Universal-only, rank-1 surface profile.\\
\Sone{} & Symmetric rank-1 profile, with scoped existential packages.\\
\Shigher{} & Full symmetric, impredicative higher-rank profile.\\
\Dep{} & The dependent-value extension of the indicated profile.\\
\Abs{} & The opaque-module abstraction discipline.\\
\Static{} & The mandatory static-checking profile, independent of rank.\\
$\Comp(T)$ & A computation with successful results in $T$, possibly failing.\\
\bottomrule
\end{tabularx}

Universal-only refers to the \emph{new explicit type-binding surface}. Ordinary
CUE instance variables retain their existential completion semantics in all three
profiles. Rank restrictions concern type quantifiers; dependence on ordinary
values is an independent design axis.

The forms \code{forall (A, B) T} and \code{exists A { ... }} bind variables
in a description. A quantified field \code{f(A, B): T} constrains one subject
at every instance. A parametric alias \code{S(A) = T} abbreviates a description
by capture-avoiding substitution; \code{S(U)} applies the alias. Instantiation
selection \code{f[U]} refers to the same subject at a universal instance.

The comparison uses proposal 4484 and its addendum at the cited repository
revision \cite{cueproposal,cueaddendum}. Stable CUE expressions retain their
existing elaboration and observation rules.

\clearpage
\pagenumbering{arabic}
\chapter{The existing constraint semantics}
\label{ch:base}
\section{Descriptions and completions}
CUE orders descriptions by subsumption and defines unification as their greatest
lower bound. At the level of successful data completions, the corresponding
set interpretation is inclusion, intersection, and union
\cite{cuespec}.

Fix a set $D$ of complete data inhabitants. A primitive description denotes a
subset of $D$; a concrete atom denotes a singleton. Write
\[
 \sem{1}=\{1\},\qquad \sem{\code{int}}=\Int,\qquad
 \sem{S\mathbin{\&}T}=\sem S\cap\sem T,\qquad
 \sem{S\mathbin{|}T}=\sem S\cup\sem T.
\]
An open record denotes the records satisfying its field predicates; a closed
record additionally constrains its permitted labels. Field absence is a
separate possibility from a present field with an unconstrained value.

An open declaration is more accurately a relation than a product of its field
marginals. For example:
\par\noindent\begin{minipage}{\linewidth}
\begin{lstlisting}
config: {
    x: int
    y: x + 1
}
\end{lstlisting}
\end{minipage}\par
has the relational interpretation
\[
 R(x,y)\equiv x\in\Int\ \land\ y=x+1.
\]
Its two fields cannot be interpreted as independent choices of integers.
Projection onto an exported interface hides the other variables existentially.
A complete exported record retains its exposed fields as the subject being
constrained.

\section{The sense in which variability is existential}
A field declaration \code{x: int} leaves a witness to be supplied by a completion.
Adding \code{x: >=0} removes negative witnesses. It imposes no obligation to
produce an implementation working uniformly for every integer.

More generally, if $\Delta$ names local witnesses and $\Gamma$ the exposed
interface, the description of that interface is
\[
 \sem{\exists\Delta\,P}
 =\{\rho\in\Env(\Gamma)\mid
       \exists\delta\in\Env(\Delta).\;P(\rho,\delta)\}.
\]
This is the existential reading of ordinary CUE variability. Built-in list and
pattern constraints may themselves validate all matching elements; a general
binder for an arbitrary semantic type is an additional facility.

Sharing precedes projection. Two contributions to one field elaborate to
$\exists x\,(P(x)\land Q(x))$. Two independently introduced witnesses elaborate
to $(\exists x\,P(x))\land(\exists y\,Q(y))$. These formulas carry different
sharing information.

\example{Independent instantiations of a description}{Existing CUE}
\par\noindent\begin{minipage}{\linewidth}
\begin{lstlisting}
#Cell: {value: int}
left:  #Cell & {value: 1}
right: #Cell & {value: 2}
\end{lstlisting}
\end{minipage}\par
Each instantiation has its own field cells. Within each cell, repeated
constraints share that cell. CUE references copy the associated expression and
rebind internal references to the corresponding copies \cite{cuerefs}.
The reference calculus therefore receives an elaborated graph with explicit
sharing, rather than equating every textual identifier with one global logic
variable.

\section{Universal implementation obligations}
A generic transformation needs one implementation whose behavior is valid for
arbitrary types. Write $\operatorname{Map}(A,B,m)$ for the requirement that $m$
applies a supplied transformation from $A$ to $B$ elementwise to a finite
$A$-list and returns the corresponding $B$-list. The two statements
\[
 \exists A,B,m.\;\operatorname{Map}(A,B,m),\qquad
 \exists m.\;\forall A,B.\;\operatorname{Map}(A,B,m)
\]
are different specifications. The first asks for some usable instance. The
second fixes an implementation and requires every instance. In particular,
existentially choosing $A=\varnothing$ can trivialize an input obligation.

This distinction applies to records as well as functions: a reusable module may
contain one operation that works at every element type, and a hidden
representation shared by several operations. The language needs to retain both
kinds of binding and their dependency order.

\section{The semantic layer of existing features}
The reference model is the \emph{completion layer}. The implementation also
retains closure provenance, scope graphs, field-presence information, and
preference information. A default-bearing description has successful
completions and a preferred subset. Default selection is an observation of that
additional structure. Comprehensions elaborate constraints using the host
language's established evaluation rules \cite{cuedefaults}.

Legacy programs retain their data observations and evaluation rules.
Refinement theorems apply to a fixed elaborated predicate and a fixed quantifier
dependency structure.

\chapter{Quantification as a set-theoretic operation}
\label{ch:quantifiers}
\section{Contexts and subjects}
\begin{definition}[Predicate fiber]
For a context $\Gamma$ with set of assignments $\Env(\Gamma)$, let
$\mathcal P_\Gamma=\Pow(\Env(\Gamma))$, ordered by inclusion. A type expression
with subject $z$ denotes a predicate in $\mathcal P_{\Gamma,z:V}$.
\end{definition}
The subject is the value described by an expression. A quantifier binds its
listed variables and leaves that subject fixed. A record field supplies a
projection of the enclosing record as the subject of its field expression.

A substitution $h:\Env(\Delta)\to\Env(\Gamma)$ induces reindexing
$h^*P=h^{-1}(P)$. It preserves all set unions and intersections.
For a fixed sort $S$, let
$\pi:\Env(\Gamma)\times S\to\Env(\Gamma)$ be the projection.
Define
\begin{align}
 \exists_\pi P
 &=\{\rho\mid\exists s\in S.\;(\rho,s)\in P\},\label{eq:exists}\\
 \forall_\pi P
 &=\{\rho\mid\forall s\in S.\;(\rho,s)\in P\}.\label{eq:forall}
\end{align}
These are predicates on the original interface, not families left outside the
semantic domain.

\begin{proposition}[Adjunctions]\label{prop:adjoints}
For predicates $P$ on the extended context and $Q$ on the base context,
\[
 \exists_\pi P\subseteq Q\ \Longleftrightarrow\ P\subseteq\pi^*Q,
 \qquad
 \pi^*Q\subseteq P\ \Longleftrightarrow\ Q\subseteq\forall_\pi P.
\]
\end{proposition}
\begin{proof}
The first equivalence expands to: every $(\rho,s)\in P$ has $\rho\in Q$.
The second expands to: for every $\rho\in Q$ and every $s\in S$,
$(\rho,s)\in P$.
\end{proof}
Thus $\exists_\pi\dashv\pi^*\dashv\forall_\pi$. This is the subset
hyperdoctrine underlying the proposal \cite{lawvere,mellies}.


\section{A single impredicative type sort}
Fix a set $V$ of operational inhabitants, including a designated set of closure
codes with captured environments. Their application relation is fixed
independently of type interpretation. Take
\[
 \U=\Pow(V),\qquad \rho(A)\in\U.
\]
The sort $\Type$ classifies predicates on $V$. It is a syntactic sort, not an
inhabitant of $V$ satisfying $\Type:\Type$. Type operators are interpreted
pointwise in a valuation:
\[
 \sem A_\rho=\rho(A),\qquad
 \sem{T\land U}_\rho=\sem T_\rho\cap\sem U_\rho,\qquad
 \sem{T\lor U}_\rho=\sem T_\rho\cup\sem U_\rho.
\]
Quantification has the same result sort:
\begin{align}
 \sem{\forall A\,T}_\rho
 &=\bigcap_{S\in\U}\sem T_{\rho[A:=S]},\label{eq:foralltype}\\
 \sem{\exists A\,T}_\rho
 &=\bigcup_{S\in\U}\sem T_{\rho[A:=S]}.\label{eq:existstype}
\end{align}
Both expressions denote subsets of $V$, hence members of $\U$. A variable may
therefore be instantiated by the interpretation of a quantified type. This is
impredicative polymorphism. Interpretation is defined by structural recursion
on the finite type expression; quantification ranges over a previously fixed
set of predicates, not over recursive evaluations of every source type.

\begin{proposition}[Impredicative substitution]\label{prop:subst}
For capture-avoiding substitution of any well-sorted type $S$, including a
quantified type,
\[
 \sem{T[S/A]}_\rho
 =\sem T_{\rho[A:=\sem S_\rho]}.
\]
\end{proposition}
\begin{proof}
Induct on $T$. Boolean and structural constructors commute with substitution.
At a fresh quantifier, the two sides take the same set-indexed intersection or
union over $\U$. The arrow case uses the fixed operational application relation.
\end{proof}

Operational arrows are predicates on a fixed set of realizers, rather than the
full set-theoretic exponentials $B^A$. This is the relevant departure from the
full-function-space interpretation obstructed by Reynolds's theorem
\cite{reynolds84}. Interpreting erased quantification by intersection has
precedents in realizability models of System F \cite{pistone}. Allowing failure
changes the arrow predicate; it does not by itself justify impredicativity.
The justification here is the fixed operational carrier and closure of
$\Pow(V)$ under the displayed operations.

A smaller Henkin family may replace $\Pow(V)$ if it is closed under every type
constructor and quantified interpretation admitted by the source fragment.
The powerset is a simple reference choice, not a representation to enumerate.
An implementation stores finite syntax and derivations. Data descriptions are
interpreted through their completion predicates; refinement of a description
is inclusion of those predicates. Type binders remain erased, so the predicate
sort cannot be reflected or used to construct a runtime self-membership test.

\section{Bounds and dependent ranges}
A type bound $A:T$ abbreviates the range
\[
 \operatorname{Sub}(T,\rho)
 =\{S\in\U\mid S\subseteq\sem T_\rho\}.
\]
A value binder $x\mathbin{\mathrm{in}}T$ ranges over $\sem T_\rho$.
Bounds may depend on earlier variables. Their context extension is the
comprehension set
\[
 \Gamma.D=\{(\rho,d)\mid\rho\in\Env(\Gamma),\ d\in D_\rho\}.
\]
Equations \eqref{eq:exists} and \eqref{eq:forall} then quantify over the fiber
$D_\rho$. This is the same construction for constant and dependent ranges.

For fixed ambient sorts, bounded quantification can be expressed by guards:
\[
 \exists a\in D\,P=\exists a\,(D(a)\land P),\qquad
 \forall a\in D\,P=\forall a\,(D(a)\Rightarrow P).
\]
The implication is a semantic operation in the Boolean predicate model; a
surface implementation can retain the guarded clause directly.

\section{Substitution and duality}
For a pullback square of context maps, reindexing commutes with both
quantifiers:
\[
 h^*(\exists_\pi P)=\exists_{\pi'}(h'^*P),\qquad
 h^*(\forall_\pi P)=\forall_{\pi'}(h'^*P).
\]
To verify either equation, fix the base assignment and compare the corresponding
fibers of the pullback. Their witnesses are in bijection. These Beck--Chevalley
laws govern capture-avoiding copying of quantified descriptions.

In each Boolean fiber,
\[
 \forall_\pi P=\neg\exists_\pi(\neg P).
\]
The semantic duality is available even in a surface fragment that exposes only
positive conjunctions, disjunctions, and universal binders.

\section{Empty ranges and uniform subjects}
The empty-index conventions are essential:
\[
 \bigcap\varnothing=V,\qquad\bigcup\varnothing=\varnothing.
\]
Hence $\forall a\in\varnothing\,P=\top$ and
$\exists a\in\varnothing\,P=\bot$. If $\U$ contains bottom, then
$\forall A\,A=\bot$. If lists are finite and bottom is admissible,
$\forall A\,\List(A)$ consists exactly of the empty list.

\example{One subject at every instance}{\Uone{}}
\par\noindent\begin{minipage}{\linewidth}
\begin{lstlisting}
empty(A): [...A]                 // []
impossible(A): A                 // _|_
impossibleRecord(n in 0 | 1): {
    value: n                     // one field cannot be both 0 and 1
}
\end{lstlisting}
\end{minipage}\par
These are intersections of predicates on one list, value, or record.
A parametric alias abbreviates a description by substitution. A
descriptor-producing function instead returns a description as a value.

\chapter{Unification, floating, and monotonicity}
\label{ch:refinement}
\section{Pointwise unification}
For a common range $D$,
\begin{equation}
 (\forall a\in D\,P(a))\land(\forall a\in D\,Q(a))
 =\forall a\in D\,(P(a)\land Q(a)).\label{eq:pointwise}
\end{equation}
This follows by regrouping intersections. Binder names are alpha-equivalent.
\example{Refining a polymorphic record across declarations}{\Uone{}}
\par\noindent\begin{minipage}{\linewidth}
\begin{lstlisting}
r(A): {wrap: func(A) -> {value: A}}
r(B): {wrap: func(B) -> {tag: "wrapped"}}

// Unified description, after alpha-renaming:
r(T): {wrap: func(T) -> {value: T, tag: "wrapped"}}
\end{lstlisting}
\end{minipage}\par
The two output obligations constrain one operation at each $T$.

Equal spelling does not identify independent existential witnesses:
\[
 (\exists a\,P(a))\land(\exists b\,Q(b))
 =\exists a,b\,(P(a)\land Q(b)).
\]
Shared existential witnesses are introduced by a shared scope and are conjoined
before their binder is discharged.

\section{Different bounds retain different obligations}
For a common body $P$,
\[
 (\forall a\in D\,P(a))\land(\forall a\in E\,P(a))
 =\forall a\in D\cup E\,P(a).
\]
Here $D\cup E$ is a union of \emph{assignments}. For type bounds,
\[
 \operatorname{Sub}(B)\cup\operatorname{Sub}(C)
 \subseteq\operatorname{Sub}(B\cup C)
\]
can be strict: mixed subtypes of $B\cup C$ need lie below neither arm.
\example{Two bounds do not become one upper bound}{\Uone{}}
\par\noindent\begin{minipage}{\linewidth}
\begin{lstlisting}
choose(A: int):    func(A, A) -> A
choose(A: string): func(A, A) -> A

// The stronger declaration also covers mixed unions:
stronger(A: int | string): func(A, A) -> A
\end{lstlisting}
\end{minipage}\par
The first pair retains both guards. With different bodies, its meaning is
$\forall A\,((A\subseteq\Int\Rightarrow P(A))\land
(A\subseteq\Str\Rightarrow Q(A)))$.

\section{Legal scope extrusion}
For $a\notin\FV(Q)$, existential Frobenius reciprocity gives
\[
 (\exists a\,P)\land Q=\exists a\,(P\land Q).
\]
For a nonempty fixed range, the universal counterpart holds:
\[
 (\forall a\,P)\land Q=\forall a\,(P\land Q).
\]
For a possibly empty dependent range, the universally quantified expression
must retain $Q$ outside the guard. Implementations store the guard and the
independent conjunct separately until inhabitation is known.

\example{An ordinary field floats as a constant obligation}{\Uone{}}
\par\noindent\begin{minipage}{\linewidth}
\begin{lstlisting}
r(A): {id: func(A) -> A}
r:    {label: "utilities"}

// The type sort contains the empty predicate:
r(A): {id: func(A) -> A, label: "utilities"}
\end{lstlisting}
\end{minipage}\par
Independent same-kind quantifiers commute. An existential and a universal
commute only when a separately justified independence law applies.
In general,
\[
 \exists y\,\forall x\,(y=x)\quad\ne\quad
 \forall x\,\exists y\,(y=x).
\]
The dependency graph records which universal choices a witness may depend on.
File order is irrelevant; dependency order is semantic.

\section{Disjunction and projection}
Existential projection preserves unions; universal projection preserves
intersections. Their opposite distributivities fail.
\example{Binder scope across a union}{\Uone{}+\Dep{}}
\par\noindent\begin{minipage}{\linewidth}
\begin{lstlisting}
a: forall (n in 0 | 1) (n | (1 - n))  // 0 | 1
b: (forall (n in 0 | 1) n) |
   (forall (n in 0 | 1) (1 - n))      // _|_
\end{lstlisting}
\end{minipage}\par
Projection of an entire record must also retain joint witnesses. Take
$R_0=\{(s,0)\}$ and $R_1=\{(s,1)\}$. Projection $p$ onto the first coordinate
satisfies
\[
 p(R_0\cap R_1)=\varnothing,
 \qquad p(R_0)\cap p(R_1)=\{s\}.
\]
Pushing record selection through a universal intersection is therefore only an
inclusion in general. Pointwise unification operates on the whole subject,
including its correlated fields.

\section{Refinement theorem}
\begin{theorem}[Monotone quantified refinement]\label{thm:monotone}
Let $P'\subseteq P$ be predicates on one elaborated context, with the same
quantifier dependency graph. Every legal composite $\mathcal Q$ of existential
and universal projections satisfies $\mathcal Q(P')\subseteq\mathcal Q(P)$.
In particular, $\mathcal Q(P\land R)\subseteq\mathcal Q(P)$.
\end{theorem}
\begin{proof}
Both quantifiers are monotone, either by \cref{prop:adjoints} or directly from
their membership clauses. Composition preserves monotonicity.
\end{proof}
A contradictory closed predicate remains contradictory. A projected atomic
result constrained to $\{v\}$ can remain $\{v\}$ or become empty. Its successful
completion cannot change to a different atom.

\section{Refinable bounds are correlated variables}
Let an existential semantic type $X$ have upper bound $\Num$, and let
\[
 P(X,f)\equiv X\subseteq\Num\ \land\ \forall A\subseteq X\,T(A,f).
\]
Adding $X\subseteq\Int$ gives $P'=P\land(X\subseteq\Int)\subseteq P$.
The result remains monotone after existentially hiding $X$.

The universal obligation for one fixed $X$ is contravariant in its bound; the
\emph{joint predicate on possible $X$ and $f$} is refined by conjunction.
A solver retains $X$ and its correlations. It does not substitute its current
upper approximation into every occurrence.

\example{A particular client and a universally capable producer}{\Uone{}}
\par\noindent\begin{minipage}{\linewidth}
\begin{lstlisting}
job: {
    x: number
    f: func(x) -> string
    result: f(x)
}
job: {x: int & >=0}
job: {x: 7, f: func(n: int) -> string: "\(n)"}
\end{lstlisting}
\end{minipage}\par
In a complete instance, \code{x} denotes one integer and the function obligation
uses its singleton type. The joint relation retains that dependency. A separate
\code{func(number) -> string} requires all numeric inputs.
\todoob The elaborator distinguishes an ordinary witness reference, whose
type-position use is a singleton, from a type-sorted variable, whose use is its
decoded semantic type. Copying a CUE description preserves that classification
and the copied witness graph.

\chapter{Function types from universal obligations}
\label{ch:functions}

\section{Values, computations, and failure}
A closure is an operational value with code, a captured environment, and a fixed
argument-binding protocol. Let $F\subseteq V$ be the set of closures. A call
packet records positional arguments, labels, and omission. A complete packet
contains values; a symbolic packet denotes a relation on its possible complete
packets, with sharing retained.

Write $\Run(f,p)\Downarrow v$ for a successful complete return. Constraint
failure, domain rejection, and an allowed execution that does not return have
no successful result. A symbolic successful result describes a set $S$ of
complete returns. Bottom is the empty completion predicate; it is not an
extra successful value in every type.

For a result predicate $B$, define
\[
 H_B(f,p)\equiv
       \forall v.\;(\Run(f,p)\Downarrow v\Rightarrow v\in B).
\]
The function type is
\begin{equation}
 A\Rightarrow B
 =\{f\in F\mid\forall p\in A.\;H_B(f,p)\}.
 \label{eq:arrow}
\end{equation}
Thus $f:A\Rightarrow B$ constrains successful returns; it promises neither a
return nor success. On symbolic descriptions the analogous condition is
$S\subseteq B$, including $S=\varnothing$. Ordinary arrows admit CUE constraint
failure without an effect annotation. A recursion policy remains an independent
execution restriction.

A computation judgment $e:\Comp(T)$ means that every successful complete
result of $e$ belongs to $T$. A value judgment $v:T$ asserts actual membership.
Only the latter supplies a witness for an existential or a value for a completed
field. In particular, a failing computation can have $\Comp(\bot)$ while no
returned value inhabits $\bot$.

A proved empty successful-result predicate reduces the corresponding
computation to \code{_|_}. A required data field with that predicate makes its
containing configuration inconsistent, even before a concrete witness is
supplied. Thus \code{x: int} conjoined with \code{x: bool} reduces to bottom.
Suspension is justified by an unresolved predicate, not by retaining a predicate
already proved empty.

A closure is itself an operational value. A never-successful closure can
inhabit $A\Rightarrow\bot$, although applying it supplies no returned value.
The empty-domain identity $\varnothing\Rightarrow B=F$ also holds. Emptiness of
a result predicate, emptiness of a set of closures, and failure of a particular
call are therefore distinct judgments. Their scopes govern bottom propagation
in \cref{sec:failure-propagation}. Explicit interfaces also undergo
relevant-consistency checking: a refuted successful result on a non-refuted
stated input region is a static interface error. This condition does not change
\cref{eq:arrow}; it restricts the descriptions admitted at source boundaries.

The static judgment in \cref{ch:residual} checks the source body and calling
convention by fixed rules. An annotation creates a proof obligation. It creates
no result filter or other executable operation.

\section{Variance and intersections}
\begin{proposition}[Arrow laws]\label{prop:arrow}
For all domain and result sets:
\begin{align*}
 A\subseteq C,\ D\subseteq B
 &\Longrightarrow(C\Rightarrow D)\subseteq(A\Rightarrow B),\\
 (A\cup C)\Rightarrow B
 &=(A\Rightarrow B)\cap(C\Rightarrow B),\\
 A\Rightarrow(B\cap D)
 &=(A\Rightarrow B)\cap(A\Rightarrow D).
\end{align*}
\end{proposition}
\begin{proof}
The first two statements restrict or split the universal input obligation.
For the third, every successful result must belong to both $B$ and $D$.
Intersecting those obligations is exactly the obligation for $B\cap D$.
Failure satisfies each implication vacuously. The identities also hold for a
nondeterministic relation when every successful result is constrained.
\end{proof}
Thus $(\Int\Rightarrow\Str)\cap(\Num\Rightarrow\Str)
=\Num\Rightarrow\Str$. More generally, an intersection of arrows preserves
input/output correlations that a single arrow may lose.

\example{A correlation-preserving overloaded contract}{\Uone{}}
\par\noindent\begin{minipage}{\linewidth}
\begin{lstlisting}
convert: (func(int) -> string) & (func(string) -> int)
\end{lstlisting}
\end{minipage}\par
One implementation must meet both obligations. A body may perform explicit
case analysis. The type intersection itself performs no branch selection.
On overlapping domains, all result obligations apply.

\section{Implementations and application refinement}
For a complete closure $f$, unification with $T$ denotes $\{f\}\cap\sem T$.
A successful conformance check preserves the closure. An incompatibility makes
that intersection empty. A symbolic closure retains its captured-environment predicates. Executable use additionally requires the mandatory
static certificate defined in \cref{ch:residual}. An unproved strict conformance
judgment is blocked rather than retained as a runtime type test.

For a set of closures $X$ and input packets $Q$, define
\[
 \App(X,Q)=\{v\mid f\in X,\ p\in Q,\ \Run(f,p)\Downarrow v\}.
\]
Joint environments are retained when $X$ and $Q$ share variables.
Then
\begin{equation}
 \App(X\cap Y,Q)\subseteq\App(X,Q)\cap\App(Y,Q).\label{eq:appmeet}
\end{equation}
Membership on the left has one closure witness that works on both sides.
Equality can fail: distinct closures can agree at a particular input while
having empty singleton intersection.

\example{Capabilities and one-call constraints}{\Uone{}}
\par\noindent\begin{minipage}{\linewidth}
\begin{lstlisting}
wide:   func(number) -> string
wide:   func(x: int) -> string: "ok"  // blocked: static domain mismatch

narrow: func(int) -> string
narrow: func(x: number) -> string: "ok" // conforming implementation
\end{lstlisting}
\end{minipage}\par
The second attachment satisfies the static contravariant interface rule.
The first is blocked by that rule even though a narrow implementation that fails
on other inputs could satisfy the wider partial-correctness predicate.
Semantic membership is intentionally more permissive than source certification.
A particular invocation is constrained with existential argument cells as in
\cref{ch:refinement}.

For one declaration subject, collect all interface conjuncts before checking its
available body. An inline result annotation and an independently declared
function type contribute to the same predicate. The checker retains the body's
inferred refinements or its proof graph; a weaker previously published result
annotation does not discard that evidence. Interface checking never modifies
that body. \Cref{sec:declaration-coherence} gives the normalization rule and
its consequences for declaration regrouping.


\section{Erased polymorphic functions and unification}
A polymorphic identity constraint denotes
\[
 I=\bigcap_{A\in\U}(A\Rightarrow A).
\]
Singleton instantiation proves that a successful return on $x$ is $x$. A member
of $I$ may fail on some inputs. Thus $I$ describes partial identities.

\example{Universal instances and explicit relevant observations}{\Uone{}+\Static{}}
\begin{lstlisting}
id(A): func(A) -> A
id:    func(int) -> int       // redundant instance

quiet(A): func(A) -> A
quiet:    func(int) -> bool   // invalid explicit interface on int
\end{lstlisting}
The ordinary integer clause supplies a relevant input region. Instantiating the
universal clause at that region gives both $\Int$ and $\Bool$ as successful-result
requirements. Their intersection is empty, so the accumulated explicit interface
is rejected without a body or a call. The source condition and its finite
instantiation rule are defined in \cref{sec:relevant-consistency}.

Denotationally the conjunction remains inhabited. For example,
\[
 q=\lambda x.\,(x\mathbin{\&}\Bool)
 \quad\Longrightarrow\quad
 q\in I\cap(\Int\Rightarrow\Bool),\qquad q(\mathsf{true})=\mathsf{true}.
\]
Behavior outside the stated integer region cannot repair its refuted successful
observations. Relevant consistency diagnoses the explicit specification; it
does not assert that the set of partial identities is empty. The predicate
$I\cap(\Int\Rightarrow\Bool)$ remains available to semantic reasoning and to
internal descriptions of failing applications.

\example{An empty value hypothesis}{\Uone{}}
\begin{lstlisting}
x: int
x: bool                      // _|_, without a supplied value
uninhabited(A): A             // _|_, by instantiation at bottom
\end{lstlisting}
These predicates have no value inhabitants. They differ from an inhabited set
of closures whose calls fail on a specified domain.

The meet operator is a refinement-preserving primitive. If $d$ and $e$ denote
completion predicates $D$ and $E$, their successful unification denotes
$D\cap E$. Hence
\[
 D\subseteq A,\ E\subseteq B
 \quad\Longrightarrow\quad D\cap E\subseteq A\cap B.
\]
No compatibility or inhabitance premise is needed. This gives the primitive
partial type $\forall A,B.\;(A\times B)\Rightarrow(A\cap B)$.
Symbolic operands retain their joint constraints; a record expression with
additional permitted fields is not identified with an exact ground record.

\example{A language operator with a polymorphic type}{\Uone{}+\Static{}}
\begin{lstlisting}
f(A): func(xs: [A, A]) -> A: xs[0] & xs[1]
x: f([{a: 1}, {b: 1}])       // {a: 1, b: 1}
y: f([1, 2])                 // _|_: conflicting concrete operands

meet(A, B): func(x: A, y: B) -> (A & B): x & y
z: meet({a: 1}, {b: true})   // {a: 1, b: true}
\end{lstlisting}
The body of \code{f} synthesizes $A\cap A=A$, so certification is local.
For \code{x}, instantiate $A$ at \code{{a: 1} | {b: 1}}. The returned record
refines both open descriptions and therefore their union. Its more precise
fields follow from evaluating the body, not from a claim that an arbitrary
inhabitant of the signature always performs a meet. The incompatible call
\code{y} is refuted when its concrete operands expose the conflict. The explicit
result $A\cap B$ of \code{meet} is not refuted under its rigid declaration
variables. It therefore passes relevant consistency. A generated instance such
as $\Int\cap\Bool$ may describe computation failure without creating a new
explicit interface obligation.

\section{Effects and conditional guarantees}
Ordinary arrows already allow constraint failure. An optional effect row records
observable foreign errors or external actions beyond that failure. For an error
set $\epsilon$, define
\[
 A\Rightarrow_\epsilon B
 =\{f\in A\Rightarrow B\mid
       \forall p\in A,e.\;(\Run(f,p)=\Raise(e)\Rightarrow e\in\epsilon)\}.
\]
This is still a partial-correctness contract. A bridge can expose
\code{func(int) -> int} when conversion failure becomes bottom, or
\code{func(int) -> int !bridge} when its error tags are observable.
Neither signature needs native machine widths. Registration must still prove
logical parameter, result, and protocol compatibility in the strict kernel.
An integer adapter cannot be certified as a number adapter merely by rejecting
fractions at runtime.

Termination, guaranteed success, and absence of a selected effect are stronger
properties with separate proof obligations. A dependent law that actually
executes a call as a required field can impose such an obligation locally;
an ordinary arrow never supplies it implicitly.

\chapter{Surface binding and conservative extension}
\label{ch:surface}
\section{Quantifier expressions}
The core surface expressions are
\par\noindent\begin{minipage}{\linewidth}
\begin{lstlisting}
forall (A, B) T
exists A { ... }
forall (A: number) func(A, A) -> [A, A]
forall (n in int & >=0) T
\end{lstlisting}
\end{minipage}\par
Parentheses group multiple binders or a constrained binder; a single simple
binder may omit them. The body is the following expression. A record body is
already delimited by braces. Quantifiers have low expression precedence, and
parentheses delimit narrower scopes inside unions and arrows.

Unadorned $A$ has sort $\Type$. \code{A: T} bounds that semantic type by
inclusion. \code{x in T} binds an ordinary value. The explicit form
\code{A in Type: T} records the type sort and its upper bound.
All binders are lexical and alpha-renamable. Bounds may reference earlier
binders and outer variables.

\section{Quantified field declarations}
The declaration
\par\noindent\begin{minipage}{\linewidth}
\begin{lstlisting}
map(A, B): func(func(A) -> B, [...A]) -> [...B]
\end{lstlisting}
\end{minipage}\par
elaborates exactly to
\par\noindent\begin{minipage}{\linewidth}
\begin{lstlisting}
map: forall (A, B) func(func(A) -> B, [...A]) -> [...B]
\end{lstlisting}
\end{minipage}\par
The field value is fixed outside the quantifier. The outer record's field
selection is retained as the subject of the quantified predicate.

For a quantified record:
\par\noindent\begin{minipage}{\linewidth}
\begin{lstlisting}
r(A): {
    id:   func(A) -> A
    wrap: func(A) -> {value: A}
}
\end{lstlisting}
\end{minipage}\par
the order is $\exists r\,\forall A\,P(r,A)$ when the surrounding record is
closed existentially. Both projections refer to the same record for all $A$.
The field-sugar form is concise for rank-1 declarations; RHS syntax also permits
nested and higher-rank positions.

\section{Quantifier blocks}
A lexical binder prefix can be written in a record block:
\par\noindent\begin{minipage}{\linewidth}
\begin{lstlisting}
x: {
    exists T: U
    forall V: W
    // The whole following declaration body is in this prefix.
    value: T
    transform: func(V) -> V
}
\end{lstlisting}
\end{minipage}\par
This elaborates to
\par\noindent\begin{minipage}{\linewidth}
\begin{lstlisting}
x: exists (T: U) forall (V: W) {
    value: T
    transform: func(V) -> V
}
\end{lstlisting}
\end{minipage}\par
The ordered prefix records dependency. The declarations after the prefix are
conjoined without source-order dependence. Contributions from other files
unify with the entire quantified subject; they do not capture these private
binder names merely by spelling them alike. A named interface may explicitly
export a sharing equation when such sharing is intended.

A syntactically asymmetric alternative requires external parentheses for
universals and block prefixes for existentials:
\par\noindent\begin{minipage}{\linewidth}
\begin{lstlisting}
x(V: W): {
    exists T: U
    value: T
    transform: func(V) -> V
}
\end{lstlisting}
\end{minipage}\par
Its prefix is $\forall V\,\exists T$, rather than $\exists T\,\forall V$.
The two spellings express different dependence. Both quantifiers still use the
dual semantic operations of \cref{ch:quantifiers}.

\section{Function-local generics and parametric aliases}
A compact function-only spelling is possible:
\par\noindent\begin{minipage}{\linewidth}
\begin{lstlisting}
map: func<A, B>(func(A) -> B, [...A]) -> [...B]
\end{lstlisting}
\end{minipage}\par
It expands to an outer universal on the function subject.

A parametric alias binds parameters in a description:
\par\noindent\begin{minipage}{\linewidth}
\begin{lstlisting}
Box(A) = {value: A}
Pair(A, B) = [A, B]
\end{lstlisting}
\end{minipage}\par
The \code{=} form defines an alias; the \code{:} form constrains a field.
Alias application substitutes its arguments capture-avoidingly into the body,
so \code{Box(int)} expands to \code{{value: int}}. Alias parameters use the
same sort and bound syntax as quantified binders; arguments must satisfy those
bounds. A nonparametric alias uses the existing form \code{let S = T}.

In contrast, \code{box(A): {value: A}} requires one value lying in every $A$
and is bottom when bottom is admissible. Reified descriptor functions retain
their descriptor arguments as correlated variables.

\code{foo[T]} selects the instance of a universally constrained subject
justified by substituting \code{T} for its binder. For a universal description,
selection refers to the same inhabitant at that instance; it retains the subject
and its universal constraints. Thus \code{id[int](3)} uses the integer instance
of \code{id}. Selection is erased. Implicit instantiation uses fresh flexible
variables at the call site and retains the universal declaration.

\section{A conservative embedding}
\begin{definition}[Legacy-data observation]
A legacy-data observation is evaluation and export using the pre-extension
operators and data formats, applied to a program in the legacy syntax. Its
completion domain is $D$, embedded into the extended $V$ by $i:D\hookrightarrow V$.
\end{definition}

\Needspace{10\baselineskip}
\begin{theorem}[Relative conservativity]\label{thm:conservative}
Suppose the extension retains the legacy elaborator, base evaluation rules,
default-selection rules, and data constructors. For every legacy description
$T$ and legacy assignment $\rho$,
\[
 i^{-1}(\sem T^{+}_{i\rho})=\sem T^{0}_{\rho}.
\]
Closed legacy programs have the same legacy-data observations under the extended
evaluator.
\end{theorem}
\begin{proof}
For data predicates, the claim is the defining restriction of each extended
primitive to $D$. Inverse image preserves intersections and unions. Record and
list constructors preserve the embedding and the field-sharing elaboration.
Existing local projection and copying operations are unchanged. The operational
statement follows by induction on the unchanged evaluation derivation, with
identical default selection and export.
\end{proof}
The theorem is observational and data-relative. Extended top additionally
admits new kinds of inhabitants in new-language contexts. This distinguishes
conservativity from equality of global carriers.

\decision{Gate new binder and arrow forms by language version or experiment
attributes. Treat \code{forall}, \code{exists}, and \code{func} as contextual
keywords in the selected syntax. Existing data files retain their parsing,
closing, copying, and evaluation behavior.}

Existing struct-as-function code remains a relational program:
\par\noindent\begin{minipage}{\linewidth}
\begin{lstlisting}
#Add: {a: int, b: int, out: a + b}
sum: (#Add & {a: 1, b: 2}).out
\end{lstlisting}
\end{minipage}\par
A function declaration adds a reusable universal obligation and a concise call
form. Refactoring to that declaration is an explicit program change, checked
against the resulting contract.

\chapter{Three implementation profiles}
\label{ch:profiles}
\section{Rank, instantiation, and interface boundaries}
Type rank measures the placement of polymorphic binders relative to function
arguments. A callback required to work at all types is higher-rank; an ordinary
higher-order function such as \code{map} can have a rank-1 type. The semantic
sort $\Type$ is impredicative. The rank-1 source profiles limit inferred and
explicit instance types to quantifier-free bodies at executable boundaries;
\Shigher{} permits explicitly supplied quantified instances. This is a source
formation restriction, not a hierarchy of semantic universes.

For the restricted profiles, use a \emph{boundary-prenex} discipline. Within one
executable signature clause, type quantifiers occur in an outer prefix over a
quantifier-free body, after justified normalization. Module interfaces are
explicit additional boundaries. Finite intersections retain several prefixed
clauses when their bounds or dependencies differ. An alias expands for the rank
check; naming a higher-rank type does not change its rank.

A positive universal may sometimes float through a result arrow when its
variable is absent from the input and the established deterministic arrow law
justifies the move. Quantifiers over argument types retain their position.
Existentials whose witnesses depend on arguments likewise retain that dependence.

\section{Full symmetric, higher-rank CUE: \Shigher{}}
\Shigher{} admits both quantifiers at arbitrary well-sorted type positions, with the
single impredicative type sort. It includes first-class existential packages,
polymorphic arguments, and nested quantified records.
\par\noindent\begin{minipage}{\linewidth}
\begin{lstlisting}
use: func(forall A func(A) -> A) -> [int, string]

#Showable: exists A {
    value: A
    show:  func(A) -> string
}
show: func(#Showable) -> string
\end{lstlisting}
\end{minipage}\par
The checker uses explicit introduction and elimination rules, scope levels, and
annotations. Higher-rank bidirectional checking provides a useful proof-theoretic
organization \cite{dunfield}, while Boolean connectives and dependent predicates
require their own sound constraint procedures.


\example{An explicitly impredicative instance}{\Shigher{}+\Static{}}
\begin{lstlisting}
let Identity = forall A func(A) -> A
id(A): func(x: A) -> A: x
again: id[Identity](id)
a: again[int](3)
b: again[string]("three")
\end{lstlisting}
The instance argument is itself universally quantified. Checking substitutes
finite type syntax; it does not search for an impredicative instance or evaluate
any type representation at runtime. Implicit higher-rank inference may request
this annotation. The mandatory checking profile in \cref{ch:residual} is
orthogonal to the three formation profiles.

\section{Symmetric, rank-1 CUE: \Sone{}}
\Sone{} admits existential and universal prefixes at declaration and module-interface
boundaries. An executable callback parameter is monomorphic at each call.
Abstract modules are opened at their module boundary, exposing a rigid local
representation name and monomorphic operations to their clients. Passing an
inline quantified package to a function, or receiving a universally quantified
callback, uses \Shigher{}.

A representative interface is
\par\noindent\begin{minipage}{\linewidth}
\begin{lstlisting}
Counter: exists State {
    zero: State
    next: func(State) -> State
    read: func(State) -> int
}
\end{lstlisting}
\end{minipage}\par
Module-level client code opens \code{Counter}, then passes \code{next} and
\code{read} as ordinary monomorphic functions. The prefix can contain both kinds
of binder, with explicit dependency order. Symmetry concerns their availability
and semantic laws at the same boundaries; it does not authorize commuting them.

This profile supports ML-style abstract module reuse with a boundary-prenex
rank restriction.

\section{Universal-only, rank-1 CUE: \Uone{}}
\Uone{} adds only explicit outer universal type binders. Existing instance variables,
local field cells, and call-instantiation variables remain existentially
interpreted. The surface need not expose existential type packaging to express
identity, mapping, composition, filtering, folds supplied as primitives, and
many generic adapter interfaces.

The same semantics still defines existential projection internally. The
implementation simply exposes a smaller formation and elimination interface.
It therefore gains universal capability obligations and compositional
polymorphic subtyping without requiring an opaque-module boundary in its first
release.

\par\noindent\begin{minipage}{\linewidth}
\begin{lstlisting}
id(A): func(A) -> A
compose(A, B, C): func(func(B) -> C, func(A) -> B) ->
    func(A) -> C
map(A, B): func(func(A) -> B, [...A]) -> [...B]
\end{lstlisting}
\end{minipage}\par
These declarations require rank-1 universals. They already distinguish reusable
implementations from existentially constrained calls.

\section{Coverage of host-language interfaces}
Rank-1 signatures describe ordinary generic APIs. Higher-rank callbacks,
abstract packages, and some lifetime-polymorphic or dependent APIs require a
richer representation.
GHC's documented distinction between rank-1 schemes and polymorphic parameters
is a direct example \cite{ghcrank}. A host binding must either express its actual
quantifier structure or export a deliberately weaker interface with that loss
recorded.

\begin{tabularx}{\textwidth}{@{}Yccc@{}}
\toprule
Facility & \Uone{} & \Sone{} & \Shigher{}\\\midrule
Outer polymorphic arrows and records & yes & yes & yes\\
Finite intersections of scheme clauses & yes & yes & yes\\
Existential instance variables & yes & yes & yes\\
Explicit abstract module type witnesses & --- & yes & yes\\
Polymorphic callback parameters & --- & --- & yes\\
First-class nested quantified packages & --- & --- & yes\\
Value dependence as a separate extension & \Dep{} & \Dep{} & \Dep{}\\
\bottomrule
\end{tabularx}

\section{Implementation choice}
\Uone{} supports rank-1 generic interfaces. \Sone{} adds explicit modular abstraction
at stable boundaries. \Shigher{} admits quantified descriptions at arbitrary
well-sorted positions. The inclusion of source fragments is conservative on
programs already admitted by the smaller profile.
A rank restriction does not by itself decide arbitrary semantic subtyping,
quantified arithmetic, or program equivalence; each checker specifies its
supported predicate fragment.

\chapter{Programs and calling conventions}
\label{ch:programs}
\section{Basic polymorphic constructions}
\example{Identity at unrelated instances}{\Uone{}}
\par\noindent\begin{minipage}{\linewidth}
\begin{lstlisting}
id(A): func(x: A) -> A: x
integer: id(7)                 // 7
text:    id("seven")           // "seven"
record:  id({port: 443})       // {port: 443}
chosen:  id[int & >=0](7)      // a checked explicit instance
\end{lstlisting}
\end{minipage}\par
One closure meets every instance. Explicit instantiation does not construct a
wrapper or remove other valid instances.

\example{Pairing preserves an inferred common refinement}{\Uone{}}
\par\noindent\begin{minipage}{\linewidth}
\begin{lstlisting}
pair(A: number): func(x: A, y: A) -> [A, A]: [x, y]
p: pair[int & >=1](2, 5)       // [2, 5]
q: pair(1, 2.5)                // [1, 2.5]
\end{lstlisting}
\end{minipage}\par
The mixed call has the admissible instance $A=\{1,2.5\}$. A shared type variable
preserves a common refinement; it does not enforce equal primitive tags.
A constant result $[0,0]$ would satisfy the unrefined numeric result type but
would fail the instance $A=\{2,5\}$. Matching primitive categories can instead
be specified by the input relation $[\Int,\Int]\cup[\mathsf{Float},\mathsf{Float}]$.

\example{Selecting a value preserves its subtype}{\Uone{}}
\par\noindent\begin{minipage}{\linewidth}
\begin{lstlisting}
min(A: number): func(x: A, y: A) -> A: {
    if x <= y {out: x}
    if x > y  {out: y}
}.out

m: min[int & >=10](12, 20)     // 12
\end{lstlisting}
\end{minipage}\par
Either branch returns an input inhabitant. By contrast, the proposed contract
\code{forall (A: number) func(A) -> A} does not justify adding one: $A=\{0\}$
is a counterexample. Numeric bounds permit numeric operations; result
refinements require their own proof.

\example{Two type parameters preserve asymmetric information}{\Uone{}}
\par\noindent\begin{minipage}{\linewidth}
\begin{lstlisting}
swap(A, B): func(p: [A, B]) -> [B, A]: [p[1], p[0]]
s: swap([7, "seven"])         // ["seven", 7]
\end{lstlisting}
\end{minipage}\par
The two coordinates retain their distinct type variables, rather than being
merged into one union.

\Needspace{180pt}
\section{Generic collection programs}
\example{Map with a single implementation}{\Uone{}}
\par\noindent\begin{minipage}{\linewidth}
\begin{lstlisting}
map(A, B): func(g: func(A) -> B, xs: [...A]) -> [...B]: [
    for x in xs {g(x)}
]
inc:  func(x: int) -> int: x + 1
text: func(x: int) -> string: "\(x)"

numbers: map(inc, [1, 2, 3])    // [2, 3, 4]
strings: map(text, [1, 2, 3])   // ["1", "2", "3"]
\end{lstlisting}
\end{minipage}\par
The body is checked with rigid $A,B$. Each call creates flexible instance
variables. Finite sets of admissible instantiations can contribute intersected
typings, as in the CDuce polymorphic application discipline \cite{cducetutorial}.

\example{Filtering and partitioning preserve elements}{\Uone{}}
\par\noindent\begin{minipage}{\linewidth}
\begin{lstlisting}
filter(A): func(p: func(A) -> bool, xs: [...A]) -> [...A]: [
    for x in xs if p(x) {x}
]
partition(A): func(p: func(A) -> bool, xs: [...A]) -> {
    yes: [...A]
    no:  [...A]
}: {
    yes: [for x in xs if p(x) {x}]
    no:  [for x in xs if !p(x) {x}]
}
positive: func(x: int) -> bool: x > 0
kept: filter(positive, [-2, 0, 4]) // [4]
\end{lstlisting}
\end{minipage}\par
The callback is pure. An effectful callback requires an explicit effect policy,
including the meaning of repeated evaluation.

\example{Flattening a generic transformation}{\Uone{}}
\par\noindent\begin{minipage}{\linewidth}
\begin{lstlisting}
flatMap(A, B): func(g: func(A) -> [...B], xs: [...A]) -> [...B]: [
    for x in xs for y in g(x) {y}
]
duplicate(A): func(x: A) -> [A, A]: [x, x]
twice: flatMap(duplicate, [1, 2]) // [1, 1, 2, 2]
\end{lstlisting}
\end{minipage}\par
\example{Mapping named record coordinates}{\Uone{}}
\par\noindent\begin{minipage}{\linewidth}
\begin{lstlisting}
mapPoint(A, B): func(g: func(A) -> B, p: {x: A, y: A}) -> {
    x: B
    y: B
}: {
    x: g(p.x)
    y: g(p.y)
}
p: mapPoint(text, {x: 2, y: 3}) // {x: "2", y: "3"}
\end{lstlisting}
\end{minipage}\par
Open/closed record constraints retain their existing meaning. A separate row
parameter can preserve additional fields when a row-polymorphic fragment is
selected.

\section{Higher-order composition and return values}
\example{Composition is rank-1}{\Uone{}}
\par\noindent\begin{minipage}{\linewidth}
\begin{lstlisting}
compose(A, B, C): func(g: func(B) -> C, f: func(A) -> B) ->
    func(A) -> C:
    func(x: A) -> C: g(f(x))

pipeline: compose(text, inc)
answer: pipeline(4)           // "5"
\end{lstlisting}
\end{minipage}\par
The returned closure captures $f$ and $g$. Its generic variables are rigid
assumptions while checking the body; their run-time representations are erased.

\example{A value-dependent closure}{\Uone{}}
\par\noindent\begin{minipage}{\linewidth}
\begin{lstlisting}
withPrefix: func(prefix: string) -> func(string) -> string:
    func(s: string) -> string: prefix + s

warn: withPrefix("warning: ")
message: warn("retry")        // "warning: retry"
\end{lstlisting}
\end{minipage}\par
Two closures from the same body with different concrete prefixes have different
captured environments.

\example{A polymorphic callback}{\Shigher{}}
\par\noindent\begin{minipage}{\linewidth}
\begin{lstlisting}
useBoth: func(p: forall A func(A) -> A) -> [int, string]: [
    p(7),
    p("seven"),
]
ok: useBoth(id)              // [7, "seven"]
onlyInt: func(x: int) -> int: x
bad: useBoth(onlyInt)         // incompatible universal obligation
\end{lstlisting}
\end{minipage}\par
The callback itself is polymorphic. A new instantiation of the surrounding
function for each use of $p$ would not prove this contract.

\example{A nested universal result}{\Shigher{} syntax; prenex normalization may apply}
\par\noindent\begin{minipage}{\linewidth}
\begin{lstlisting}
constant(A): func(x: A) -> (forall B func(B) -> A):
    func(_) -> A: x

seven: constant(7)
a: seven("ignored")          // 7
b: seven(true)               // 7
\end{lstlisting}
\end{minipage}\par
The resulting closure depends on $x$ and works uniformly for every subsequent
argument type. Moving $B$ outward is permitted only through the proven positive
arrow law, with $B$ absent from the domain.

\section{Quantified records and independent refinement}
\example{A reusable polymorphic utility record}{\Uone{}}
\par\noindent\begin{minipage}{\linewidth}
\begin{lstlisting}
utilities(A): {
    id:   func(x: A) -> A: x
    wrap: func(x: A) -> {value: A}: {value: x}
}
utilities: {name: "core"}

wrapped: utilities.wrap("payload") // {value: "payload"}
\end{lstlisting}
\end{minipage}\par
The \code{name} field is a constant conjunct on the same record.

\example{A cross-file behavioral strengthening}{\Uone{}}
\par\noindent\begin{minipage}{\linewidth}
\begin{lstlisting}
// interface.cue
wrap(A): func(A) -> {value: A}

// metadata.cue
wrap(A): func(A) -> {tag: "wrapped"}

// implementation.cue
wrap(A): func(x: A) -> {value: A, tag: "wrapped"}: {
    value: x
    tag: "wrapped"
}
\end{lstlisting}
\end{minipage}\par
The body is checked against both pointwise result constraints. No declaration
replaces or takes precedence over another.

\example{An intersection of domain-specific guarantees}{\Uone{}}
\par\noindent\begin{minipage}{\linewidth}
\begin{lstlisting}
classify: func(x: int) -> ("negative" | "nonnegative"): {
    if x < 0  {out: "negative"}
    if x >= 0 {out: "nonnegative"}
}.out
classify: (func(int & <0) -> "negative") &
          (func(int & >=0) -> "nonnegative")

c: classify(-3)              // "negative"
\end{lstlisting}
\end{minipage}\par
An ordinary union of function values describes alternatives, whereas this
intersection describes one value satisfying both implications.

\example{Unresolved alternatives and defaults}{\Uone{}}
\par\noindent\begin{minipage}{\linewidth}
\begin{lstlisting}
f: (func(x: int) -> int: x) |
   (func(x: int) -> int: x + 1)
a: f(3)                     // incomplete operational choice

preferred: *(func(x: int) -> int: x) |
            (func(x: int) -> int: x + 1)
b: preferred(3)             // 3, under legacy default selection
\end{lstlisting}
\end{minipage}\par
Constraint propagation can retain the set of alternative completions. Execution
uses the ordinary concreteness/default discipline to select a callable value.

\section{Argument forms and call packets}
The five parameter forms in proposal 4484 all elaborate to packet predicates
\cite{cueproposal}. Named parameters admit a positional packet or the matching
label; required and optional name-only parameters use labels; anonymous and
aliased positional parameters use positions only. Packet binding precedes value
constraint solving.

\example{Positional, named, and mixed calls}{\Uone{}}
\par\noindent\begin{minipage}{\linewidth}
\begin{lstlisting}
sum: func(a: int, b: int) -> int: a + b
p: sum(2, 3)                // 5
q: sum(a: 2, b: 3)          // 5
r: sum(2, b: 3)             // 5
\end{lstlisting}
\end{minipage}\par

\example{Name-only required and optional parameters}{\Uone{}}
\par\noindent\begin{minipage}{\linewidth}
\begin{lstlisting}
key: func(a!: int, note?: string) -> int: a
x: key(a: 4)                // 4; note is absent
z: key(a: 4, note: "audit") // 4
bad: key(4)                 // invalid packet: a is name-only
\end{lstlisting}
\end{minipage}\par
A body that reads an absent optional parameter must prove presence or produce a
checked absence failure. Optionality is an absence alternative, not an unknown
present value.

\example{Positional-only aliases and anonymous parameters}{\Uone{}}
\par\noindent\begin{minipage}{\linewidth}
\begin{lstlisting}
scale: func(_~value: int, _: int) -> int: value * 2
x: scale(5, 99)             // 10
bad: scale(value: 5, 99)     // invalid label/order
\end{lstlisting}
\end{minipage}\par
The positional alias names a local body variable without adding a public
label. Renaming it preserves the packet protocol.

\example{Hidden labels and attributes}{\Uone{}}
\par\noindent\begin{minipage}{\linewidth}
\begin{lstlisting}
package internal
measure: func(_value: int @go(Int), unit!: string) -> int: _value
inside: measure(_value: 7, unit: "ms")
// Another package can use the positional slot, not the hidden label.
\end{lstlisting}
\end{minipage}\par
Hidden labels include their package identity. Attributes are retained as
metadata; only a declared extension assigns them operational significance.

\example{Errors distinguish presence from values}{\Uone{}}
\par\noindent\begin{minipage}{\linewidth}
\begin{lstlisting}
a: sum(1, a: 1)            // duplicate binding, even with equal values
b: sum(1, 2, 3)            // surplus positional argument
c: sum(a: 1, b: 2, c: 3)    // unknown label
d: sum(1)                  // required argument missing
e: sum(_, 2)               // supplied but incomplete, not missing
\end{lstlisting}
\end{minipage}\par
A duplicate packet is invalid before field unification. Conjoining equal values
does not erase a duplicate-binding error.

\section{Defaults}
\example{Omission defaults and caller preferences}{\Uone{}}
\par\noindent\begin{minipage}{\linewidth}
\begin{lstlisting}
add: func(a: int, b: int = 10) -> int: a + b
x: int | *5
omitted:  add(1)            // 11
supplied: add(1, x)         // 6
unknown:  add(1, _)         // incomplete; b was supplied
\end{lstlisting}
\end{minipage}\par
An omission default belongs to the closure protocol. A supplied argument bypasses
it and retains its own ordinary preference information. The default expression
is checked against its parameter description in the declaration environment.

\example{A defaulted name-only argument}{\Uone{}}
\par\noindent\begin{minipage}{\linewidth}
\begin{lstlisting}
render: func(text: string, suffix!: string = "!") -> string:
    text + suffix
x: render("hello")                 // "hello!"
y: render("hello", suffix: "?")   // "hello?"
\end{lstlisting}
\end{minipage}\par
Defaulted required-label parameters may be omitted, while optional parameters
retain absence semantics and have no omission default.

\example{A configurable default is an explicit adapter}{\Uone{}}
\par\noindent\begin{minipage}{\linewidth}
\begin{lstlisting}
raw: func(x: int) -> int: x
configured: func(x: int = 2) -> int: raw(x)
answer: configured()        // 2
\end{lstlisting}
\end{minipage}\par
The adapter supplies the protocol that accepts the empty packet. A contract
mentioning that protocol requires an implementation possessing it. It does not
mutate a previously supplied closure.

\section{Open interfaces and partial application}
An open signature describes a partial \emph{protocol interface}. Elaborate its
ellipsis to an existential row of further parameter declarations, retaining
that witness with the function. Its callable packet domain is determined only
when the required row is known. This is different from promising that every
completion of the row can be omitted. \Uone{} retains this row as an internal
existential witness, just as it retains call-instantiation variables; explicit
existential type packaging is a separate surface facility.

\example{An open interface completed by an implementation}{\Uone{}, implicit protocol row}
\par\noindent\begin{minipage}{\linewidth}
\begin{lstlisting}
T: func(a: int, ...) -> number
f: T & (func(a: int, b: int) -> int: a + b)
x: f(1, 2)                 // 3
y: f(1)                    // missing b
\end{lstlisting}
\end{minipage}\par
The row witness is completed by \code{b: int}. A closed one-argument capability
admits an implementation with additional omittable parameters because its
promised one-argument packets remain covered.

\example{Chained partial application}{\Uone{}}
\par\noindent\begin{minipage}{\linewidth}
\begin{lstlisting}
add3: func(a: int, b: int, c: int) -> int: a + b + c
f: add3(1, ...)
g: f(2, ...)
x: g(3)                    // 6
y: add3(a: 1, ...)(2, 3)    // 6
\end{lstlisting}
\end{minipage}\par
A partial closure stores bindings to original parameter slots. Supplied values
are checked immediately; the body runs on completion. An unbound default is
preserved and a bound argument suppresses that default.

\example{A residual contract on a partial closure}{\Uone{}}
\par\noindent\begin{minipage}{\linewidth}
\begin{lstlisting}
f: add3(1, ...)
f: func(int, int) -> int
x: f(2, 3)                 // 6
\end{lstlisting}
\end{minipage}\par
The residual protocol is a normal callable contract. Normalizing the original
slot bindings makes partial applications comparable independently of whether
arguments were supplied by label or position.

\section{Identity, captured values, and result combination}
Use a closure descriptor $(o,\gamma,\beta)$: declaration origin $o$, captured
free environment $\gamma$, and normalized partial bindings $\beta$.
The origin survives copying; irrelevant ambient fields are absent from
$\gamma$. Symbolic captures remain constraint variables.

\example{Self-unification after copying and embedding}{\Uone{}}
\par\noindent\begin{minipage}{\linewidth}
\begin{lstlisting}
x: {f: func(y: int) -> int: y}
a: x.f & x.f
b: x.f & {x}.f
c: x.f & (x & {g: 1}).f
r: [a(2), b(2), c(2)]       // [2, 2, 2]
\end{lstlisting}
\end{minipage}\par

\example{Equal source, unequal captures}{\Uone{}}
\par\noindent\begin{minipage}{\linewidth}
\begin{lstlisting}
d: [for v in [1, 2] {func(y: int) -> int: y + v}]
e: d[0] & d[1]             // _|_: captures differ
\end{lstlisting}
\end{minipage}\par
For one origin, captured-environment constraints unify structurally. Distinct
known origins denote distinct intensional values. This supplies idempotence
without deciding equality of mathematical functions.

\example{Combining result descriptions explicitly}{\Uone{}}
\par\noindent\begin{minipage}{\linewidth}
\begin{lstlisting}
left:  func(x: int) -> {a: int}: {a: x}
right: func(x: int) -> {b: int}: {b: x}
both: func(x: int) -> {a: int, b: int}: left(x) & right(x)
r: both(3)                  // {a: 3, b: 3}
\end{lstlisting}
\end{minipage}\par
The body is a new implementation whose result is the CUE unification of the two
result descriptions. The primitive meet rule gives its result their intersection.
Conflicting callbacks yield bottom on the affected calls; their consistency need
not be proved when the function is defined.

\section{Validators and iteration}
\example{A predicate used as an ordinary constraint}{\Uone{}}
\par\noindent\begin{minipage}{\linewidth}
\begin{lstlisting}
positive: func(x: int) -> bool: x > 0
#Positive: {value: int, _valid: true & positive(value)}
n: (#Positive & {value: 7}).value   // 7
bad: #Positive & {value: -1}       // _|_
\end{lstlisting}
\end{minipage}\par
A validator is a predicate on an existential subject. It can be expressed with
existing struct machinery once predicate functions are available.

\example{Portable foreign contract}{\Uone{} with checked effects}
\par\noindent\begin{minipage}{\linewidth}
\begin{lstlisting}
externF: extern func(int) -> int !bridge
// A native-range failure is a bridge outcome.
// A fractional input is outside the logical domain.
\end{lstlisting}
\end{minipage}\par
The bridge checks logical kind, representation, conversion, and returned
values as specified by the source proposal \cite{cueaddendum}. The effect
annotation records those failures; it introduces no additional conversion
stage. Tool-level effects execute in an explicit sequencing context, while
constraint unification operates on their descriptions.

\example{Nested calls remain ordinary evaluation}{\Uone{}}
\par\noindent\begin{minipage}{\linewidth}
\begin{lstlisting}
twice: func(n: int) -> int: n + n
value: twice(twice(twice(2))) // 16
\end{lstlisting}
\end{minipage}\par
The initial execution profile preserves the host cycle discipline. A structurally
recursive extension can admit calls justified by a decreasing finite measure.
This choice is independent of quantifier rank.

\example{A structurally recursive fold}{\Uone{} with structural recursion}
\par\noindent\begin{minipage}{\linewidth}
\begin{lstlisting}
fold(A, B): func(step: func(B, A) -> B, seed: B, xs: [...A]) -> B: {
    if len(xs) == 0 {out: seed}
    if len(xs) > 0 {out: fold(step, step(seed, xs[0]), xs[1:])}
}.out
plus: func(x: int, y: int) -> int: x + y
sum: fold(plus, 0, [1, 2, 3]) // 6
\end{lstlisting}
\end{minipage}\par
Each recursive call reduces list length. The termination proof is independent
of polymorphism; a nonrecursive profile may supply the same contract as a trusted
primitive.

\chapter{Dependent descriptions and functions}
\label{ch:dependent}
\section{Dependent contexts}
A dependent context extends $\Gamma$ by a value $x$ satisfying a description
$A$ in that context:
\[
 \Env(\Gamma,x:A)
 =\{(\rho,x)\mid\rho\in\Env(\Gamma),\ x\in\sem A_\rho\}.
\]
The existential and universal adjoints already apply to its projection.
Dependence therefore changes the fibers over which quantification ranges,
rather than introducing another interpretation of refinement.

A dependent function contract uses its actual argument in the result predicate:
\[
 \Pi_{x:A}B(x)
 =\{f\in F\mid\forall x\in A.\;H_{B(x)}(f,x)\}.
\]
A dependent pair retains the witness:
\[
 \Sigma_{x:A}B(x)=\{(x,y)\mid x\in A,\ y\in B(x)\}.
\]
Logical $\forall x\,B(x)$ constrains \emph{one subject} at every $x$; a dependent
function instead relates inputs to results. Logical $\exists x\,B(x)$ projects
away $x$; a dependent pair retains it. The distinction is determined by the
subject and the application or pairing relation.

\decision{Within a dependent function signature, parameters enter scope from
left to right. A result constraint can refer to all parameters. Default
expressions can depend only on parameters whose binding is already established.
The dependency graph determines evaluation order; it is independent of the
source order of unrelated record declarations.}

\example{A directly dependent result}{\Uone{}+\Dep{}}
\par\noindent\begin{minipage}{\linewidth}
\begin{lstlisting}
successor: func(x: int) -> (x + 1): x + 1

// Equivalent pointwise specification:
successor(n in int): func(n) -> (n + 1)
\end{lstlisting}
\end{minipage}\par
The second form quantifies over input singletons. The first uses a dependent
result predicate on the actual argument.

\example{A dependent parameter constraint}{\Uone{}+\Dep{}}
\par\noindent\begin{minipage}{\linewidth}
\begin{lstlisting}
interval: func(lo: int, hi: int & >=lo) -> {
    low: lo
    high: hi
}: {low: lo, high: hi}

ok: interval(3, 7)          // {low: 3, high: 7}
bad: interval(7, 3)         // domain contradiction
\end{lstlisting}
\end{minipage}\par
The admitted domain is a relation on two arguments. It is not the Cartesian
product of their marginal types.

\Needspace{180pt}
\section{Indexed finite sequences}
The following aliases use the list validators supplied by CUE
\cite{cuelist}:
\par\noindent\begin{minipage}{\linewidth}
\begin{lstlisting}
import "list"

let Nat = int & >=0
Vec(A, n in Nat) = [...A] &
    list.MinItems(n) & list.MaxItems(n)
Fin(n in Nat) = int & >=0 & <n
\end{lstlisting}
\end{minipage}\par
Their reference semantics is
\[
 \sem{\operatorname{Vec}(A,n)}
   =\{xs\in\List(A)\mid |xs|=n\},\qquad
 \sem{\operatorname{Fin}(n)}=\{0,\ldots,n-1\}.
\]
A symbolic index remains a constraint variable. The checker uses the length
predicate and its lemmas; concrete validator execution is one implementation
strategy for ground indices.

\example{Length-preserving map}{\Uone{}+\Dep{}}
\par\noindent\begin{minipage}{\linewidth}
\begin{lstlisting}
map(A, B, n in Nat): func(
    g: func(A) -> B,
    xs: Vec(A, n),
) -> Vec(B, n): [for x in xs {g(x)}]
\end{lstlisting}
\end{minipage}\par
The universal $n$ records a property of the input, rather than supplying a
runtime integer. Its value is available as \code{len(xs)} when evaluation
requires it.

\example{Safe indexing, including the zero-length case}{\Uone{}+\Dep{}}
\par\noindent\begin{minipage}{\linewidth}
\begin{lstlisting}
index(A, n in Nat): func(xs: Vec(A, n), i: Fin(n)) -> A: xs[i]

x: index(["a", "b"], 1)    // "b"
bad: index([], 0)           // 0 is outside Fin(0)
\end{lstlisting}
\end{minipage}\par
For $n=0$, the packet domain is empty. The generic implementation remains valid;
there is no admitted out-of-range call.

\example{Head with an explicit nonempty precondition}{\Uone{}+\Dep{}}
\par\noindent\begin{minipage}{\linewidth}
\begin{lstlisting}
head(A, n in int & >=1): func(xs: Vec(A, n)) -> A: xs[0]
first: head([4, 5])         // 4
bad: head([])               // no admissible instance
\end{lstlisting}
\end{minipage}\par

\example{Appending adds lengths}{\Uone{}+\Dep{}}
\par\noindent\begin{minipage}{\linewidth}
\begin{lstlisting}
append(A, n in Nat, m in Nat): func(
    xs: Vec(A, n), ys: Vec(A, m),
) -> Vec(A, n + m): [for x in xs {x}, for y in ys {y}]

v: append([1, 2], [3])      // [1, 2, 3]
\end{lstlisting}
\end{minipage}\par
The proof combines the cardinalities of two concatenated finite comprehensions.

\example{Zip rejects unequal lengths}{\Uone{}+\Dep{}}
\par\noindent\begin{minipage}{\linewidth}
\begin{lstlisting}
zip(A, B, n in Nat): func(xs: Vec(A, n), ys: Vec(B, n)) ->
    Vec([A, B], n): [for i, x in xs {[x, ys[i]]}]

z: zip([1, 2], ["a", "b"]) // [[1, "a"], [2, "b"]]
bad: zip([1], ["a", "b"])   // inconsistent length witness
\end{lstlisting}
\end{minipage}\par
The same existential instance of $n$ must satisfy both input constraints at a
call. Separate instances per argument would lose this invariant.

\example{Split at a bounded position}{\Uone{}+\Dep{}}
\par\noindent\begin{minipage}{\linewidth}
\begin{lstlisting}
split(A, n in Nat, k in int & >=0 & <=n): func(
    xs: Vec(A, n), cut: k,
) -> {left: Vec(A, k), right: Vec(A, n - k)}: {
    left: xs[:cut]
    right: xs[cut:]
}
\end{lstlisting}
\end{minipage}\par
The runtime parameter \code{cut} reifies the index used by slicing. The endpoints
$k=0$ and $k=n$ are admitted.

\example{Replication reifies its count}{\Uone{}+\Dep{}}
\par\noindent\begin{minipage}{\linewidth}
\begin{lstlisting}
replicate(A, n in Nat): func(x: A, count: n) -> Vec(A, n): [
    for i in list.Range(0, count, 1) {x}
]
r: replicate("x", 3)      // ["x", "x", "x"]
\end{lstlisting}
\end{minipage}\par
An erased type or value index cannot be inspected by a body merely because it
appears in a contract. Here the actual count parameter supplies the needed
runtime information. A checked library bridge can carry its declared allocation
or representation effects; the finite reference range operation is total.

\example{A dependent result with a hidden length}{\Shigher{}+\Dep{}}
\par\noindent\begin{minipage}{\linewidth}
\begin{lstlisting}
filterBounded(A, n in Nat): func(
    p: func(A) -> bool, xs: Vec(A, n),
) -> (exists (k in int & >=0 & <=n) Vec(A, k)):
    [for x in xs if p(x) {x}]
\end{lstlisting}
\end{minipage}\par
The existential witness can depend on the input and predicate. A \Uone{}+\Dep{} surface
can express the same projected result using \code{[...A] & list.MaxItems(n)}.
A witness-bearing version returns a record \code{{length: k, values: Vec(A,k)}}.

\section{Matrix dimensions, including empty matrices}
Retain dimensions as data when empty collections do not determine them:
\par\noindent\begin{minipage}{\linewidth}
\begin{lstlisting}
Matrix(A, r in Nat, c in Nat) = {
    rows: r
    cols: c
    data: Vec(Vec(A, c), r)
}
\end{lstlisting}
\end{minipage}\par
A zero-row matrix still records its column count. This makes transposition and
composition defined at all admitted dimensions.

\example{Transpose with a precise shape}{\Uone{}+\Dep{}}
\par\noindent\begin{minipage}{\linewidth}
\begin{lstlisting}
transpose(A, r in Nat, c in Nat): func(m: Matrix(A, r, c)) ->
    Matrix(A, c, r): {
        rows: m.cols
        cols: m.rows
        data: [
            for j in list.Range(0, m.cols, 1) {
                [for row in m.data {row[j]}]
            }
        ]
    }
\end{lstlisting}
\end{minipage}\par
Every $j$ is below the length $c$ of every row. The outer length is $c$ and the
inner length is $r$, including either zero dimension.

\example{Dot product with matched dimensions}{\Uone{}+\Dep{}; finite fold primitive}
\par\noindent\begin{minipage}{\linewidth}
\begin{lstlisting}
addNumber: func(x: number, y: number) -> number: x + y

dot(n in Nat): func(xs: Vec(number, n), ys: Vec(number, n)) ->
    number: fold(addNumber, 0, [for i, x in xs {x * ys[i]}])
\end{lstlisting}
\end{minipage}\par
Products and sums return \code{number}; arbitrary input refinements are not
assumed closed under arithmetic.

\example{Dimension-safe matrix multiplication}{\Uone{}+\Dep{}}
\par\noindent\begin{minipage}{\linewidth}
\begin{lstlisting}
matmul(r in Nat, k in Nat, c in Nat): func(
    x: Matrix(number, r, k), y: Matrix(number, k, c),
) -> Matrix(number, r, c): {
    rows: x.rows
    cols: y.cols
    data: [
        for row in x.data {
            [for j in list.Range(0, y.cols, 1) {
                dot(row, [for yrow in y.data {yrow[j]}])
            }]
        }
    ]
\end{lstlisting}
\end{minipage}\par
A single instance of $k$ links the two inputs and justifies each dot product.
The output carries the surviving dimensions $r,c$.

\section{Protocol descriptions and refinement predicates}
\example{A request-indexed response}{\Uone{}+\Dep{} with a checked service}
\par\noindent\begin{minipage}{\linewidth}
\begin{lstlisting}
let Request = {kind: "lookup", key: string} |
              {kind: "size", items: [..._]}
Reply(q in Request) = {
    if q.kind == "lookup" {result: string}
    if q.kind == "size"   {result: int & >=0}
}
request: extern func(q: Request) -> Reply(q) !service
\end{lstlisting}
\end{minipage}\par
The result predicate is evaluated in the request's environment. Replacing it by
an independent union of responses would admit mismatched request/response pairs.

\example{A relation between input and output records}{\Uone{}+\Dep{}}
\par\noindent\begin{minipage}{\linewidth}
\begin{lstlisting}
annotate(A): func(x: A, tag: string) -> {
    value: x
    label: tag
}: {value: x, label: tag}
\end{lstlisting}
\end{minipage}\par
The field labels differ from the parameter names. The result predicate states
equality with the corresponding argument cells and retains the input/output
correlation through subsequent record refinement.

\example{Clamping to an argument-dependent interval}{\Uone{}+\Dep{}}
\par\noindent\begin{minipage}{\linewidth}
\begin{lstlisting}
clamp: func(lo: int, hi: int & >=lo, x: int) -> (
    int & >=lo & <=hi
): {
    if x < lo {out: lo}
    if x > hi {out: hi}
    if x >= lo && x <= hi {out: x}
}.out

bounded: clamp(0, 10, 14) // 10
\end{lstlisting}
\end{minipage}\par
The three branches partition the input domain. The dependency $hi\ge lo$
justifies each branch's interval postcondition. The result is established by
linear arithmetic and branch refinement.

\section{Proof fragments and erasure}
Dependent specifications can be verified by a hierarchy of procedures:
structural length reasoning, quantifier-free linear integer arithmetic,
finite disjunction splitting, and explicit proof certificates for stronger
predicates. Refinement-type inference provides a precedent for combining type
structure with restricted automated implication checking \cite{liquid}.

The logical reference model admits well-formed dependent predicates.
Every implementation-to-contract implication is proved before certification.
A difficult first-order predicate can instead be tested by an explicit source
assertion such as \code{x & >0}. The assertion has a local typing rule that
establishes its predicate on successful results; executing it can fail or remain
pending. \Cref{sec:dynamic-admission} defines this boundary. Predicates on
functions, abstract types, or infinite universal ranges require a static proof
or an explicitly justified source operation.

All function-type annotations, type binders, instantiation selections, and
proof arguments erase. Executable assertions remain because they are terms.
Value indices used computationally remain ordinary arguments or fields, as in
replication and matrices. \Cref{thm:erasure} gives erasure for both inline and
separately contributed annotations. Explicit refinement systems provide a
related separation of proof evidence from executable code \cite{explicitrefine}.

\chapter{Existentials and modular reuse}
\label{ch:modules}
\section{A shared hidden representation}
An existential record describes several fields using one representation witness:
\par\noindent\begin{minipage}{\linewidth}
\begin{lstlisting}
#Counter: exists State {
    zero: State
    next: func(State) -> State
    read: func(State) -> (int & >=0)
}
\end{lstlisting}
\end{minipage}\par
Its meaning is $\bigcup_{S\in\U}\sem{\operatorname{Counter}(S)}$.
One $S$ must make all three fields satisfy the interface. Replacing it with
separate existential witnesses for the three fields loses the relationships
that make composition type-correct.

This is the type-theoretic basis of abstract data types and modules identified
by Mitchell and Plotkin \cite{mitchellplotkin}. In CUE, an additional distinction
matters: existential projection is a logical operation, while \emph{information
hiding} is a restriction on the observations available at an interface.

\section{Sealing and opening}
A \emph{seal} turns an implementation of an existential interface into an opaque
public view. For an unbounded representation sort $S$, a private representation
$A$ is exposed through a fresh abstract copy
\[
 \widehat A_\kappa=\{(\kappa,a)\mid a\in A\},
\]
where $\kappa$ is a seal identity. Exported operations wrap and unwrap that copy
through authorized adapters. The public existential witness is $\widehat A_\kappa$.
Thus the set interpretation of the interface remains an existential union.

Public operations are exposed through opaque capability handles associated
with the seal and the interface's declared export/alias graph. Their private
code origins and captures are not observable through that view. Copying a sealed
module preserves its seal and exported aliasing; explicitly constructing a fresh
sealed module introduces a fresh seal.

Sealing and opening have the forms
\par\noindent\begin{minipage}{\linewidth}
\begin{lstlisting}
seal Interface with (HiddenType = Representation) {
    // implementation fields, checked with the private representation
}

open moduleValue as (AbstractType, View) {
    // AbstractType is rigid; View has the opened interface.
}
\end{lstlisting}
\end{minipage}\par
Opening introduces a local rigid type and preserves sharing across the view's
field projections. The type cannot escape except under a corresponding
existential package. \Sone{} uses this operation at module boundaries; \Shigher{} admits it
inside functions over first-class packages.

A transparent bounded existential still exposes every operation justified by its
bound. A fully opaque initial sealing profile uses unbounded representation
sorts and explicit exported operations. A bound such as \code{State: number}
would expose numeric structure and entails a stronger observation-preservation
condition on any representation change.

\section{Two implementations and one client}
\example{An integer representation}{\Sone{}+\Abs{}}
\par\noindent\begin{minipage}{\linewidth}
\begin{lstlisting}
counterV1: seal #Counter with (State = int & >=0) {
    zero: 0
    next: func(x: int & >=0) -> (int & >=0): x + 1
    read: func(x: int & >=0) -> (int & >=0): x
}
\end{lstlisting}
\end{minipage}\par

\example{A record representation}{\Sone{}+\Abs{}}
\par\noindent\begin{minipage}{\linewidth}
\begin{lstlisting}
let CounterRep = close({count: int & >=0})
counterV2: seal #Counter with (State = CounterRep) {
    zero: {count: 0}
    next: func(s: CounterRep) -> CounterRep:
        {count: s.count + 1}
    read: func(s: CounterRep) -> (int & >=0): s.count
}
\end{lstlisting}
\end{minipage}\par
The abstract carrier in each public view has no accessible \code{count} field
or numeric representation. Its structure is the interface's structure.

\example{A client checked once}{\Sone{}+\Abs{} with a rank-1 utility}
\par\noindent\begin{minipage}{\linewidth}
\begin{lstlisting}
useCounter(S): func(c: {
    zero: S
    next: func(S) -> S
    read: func(S) -> (int & >=0)
}) -> (int & >=0): c.read(c.next(c.next(c.zero)))

v1: (open counterV1 as (S, C) {out: useCounter[S](C)}).out
v2: (open counterV2 as (S, C) {out: useCounter[S](C)}).out
// Both observations are 2.
\end{lstlisting}
\end{minipage}\par
The client uses only the abstract operations. Its generic declaration is
rank-1. The hidden type is opened at the module boundary and then supplied as a
rigid instance.

\example{A module law strengthened by universals}{\Sone{}+\Dep{}+\Abs{}}
\par\noindent\begin{minipage}{\linewidth}
\begin{lstlisting}
#CounterLawful: exists State {
    zero: State
    next: func(State) -> State
    read: func(State) -> (int & >=0)
    _zeroLaw: true & (read(zero) == 0)
    forall (s in State) {
        _stepLaw: true & (read(next(s)) == read(s) + 1)
    }
}
\end{lstlisting}
\end{minipage}\par
The law is a proof obligation on the module subject, not an instruction to
enumerate all states. Both implementations discharge it by substitution into
their bodies. The existential representation scope encloses the universal
state-law scope.

\section{Unification is an observation}
Opaque values can participate in CUE unification. In the basic sealed model,
\[
 (\kappa,a)\mathbin{\&}(\kappa,b)
 =\begin{cases}(\kappa,a),&a=b,\\\bot,&a\ne b.\end{cases}
\]
Consequently, a representation-independence theorem must preserve equality and
conflict, not merely the outputs of named methods.

\begin{definition}[Equality-respecting representation simulation]
For two representations $A_1,A_2$, a simulation is a bijection
$r:A_1\to A_2$ on the entire carriers admitted by the public interface,
together with corresponding exported capability handles, such that constants
correspond, every public operation commutes with $r$ on successful results,
and failure, divergence, and public effects correspond as well. The carriers include the abstract states over which a client
may quantify, as well as those reached by executing operations.
\end{definition}
For the counters, $r(n)=\{\mathit{count}:n\}$ is such a bijection. It preserves
zero, successor, reading, and equality. The abstract public seal identities are
related consistently across the two executions.

Extend $r$ componentwise to data constructors and to sets of completions by
direct image. A bijection preserves intersections, unions, emptiness, and
singletons. Its transport of abstract type assignments induces bijections on
the corresponding admissible predicate families.

\begin{theorem}[Representation-independent replacement]\label{thm:abstraction}
Assume an equality-respecting simulation, scope-preserving sealing and opening,
and admissible type families and public operational inhabitants closed under
the induced transport. Let the client's observations be confined to the declared
public interface, ordinary data, unification, and the admitted effects. Replacing one sealed
implementation by the related implementation preserves every client data
observation and its success/conflict outcome.
\end{theorem}
\begin{proof}
Interpret the client in the two related public environments. Constants and
public operations preserve the relation by hypothesis. Reindexing and data
constructors preserve it componentwise. Direct image along a bijection commutes
with unions and intersections. Quantifiers range over corresponding assignments,
so their unions and intersections commute with transport. Application preserves
the relation by the operation simulation. Sealing preserves the public alias
structure and hides implementation descriptors. Induction over the elaborated
client derivation therefore relates its resulting completion predicates. At
ordinary exported data, the relation is equality; emptiness also agrees.
\end{proof}
The abstraction discipline is related to the classical logical-relations proof
of representation independence \cite{mitchell86,reynolds83}. The explicit
equality condition is required by CUE's observable meet operation.

A many-to-one abstraction relation requires additional treatment. If multiple
private cache states represent the same logical state, their raw equality can
be observable through unification. An interface may instead expose a certified
quotient carrier and an equality congruence. Its sealing adapter then presents
logical equivalence classes, and the theorem applies to the resulting public
carrier. The implementation must realize that equality soundly.

\section{The scope of the evolution guarantee}
A partial function interface alone does not ensure preservation of successful
calls across replacement. The outcome-preserving simulation in
\cref{thm:abstraction} is the additional condition: replacing a successful
operation by one that always fails has the same partial result type but violates
that simulation. Required law fields are also preserved as observations.

The theorem yields a checkable class of representation changes for which all
admissible clients continue to work. Its hypotheses include behavior and
observable equality. A signature that merely types \code{next} and \code{read}
does not assert that \code{read(next(s)) = read(s)+1}; that equation is a separate
interface law or simulation obligation.

For one-way evolution, an observation-preserving simulation can relate the old
interface to a designated old-compatible view of the new implementation. New
operations and states may be exposed through an extended view. If additional
states are exposed through the old abstract type itself, preservation also
requires every old universally quantified obligation to hold on those states;
agreement merely on previously executed states is insufficient. Effect
guarantees and public serialization are part of the same simulation.

Here a \emph{forward-compatible client} is a client whose old interface contract
continues to hold when a later, related implementation is substituted. This is
the usual backward-compatible replacement direction for the library. The
quantified interface and simulation state the precise class of permitted future
changes.

Ordinary exposed structural schemas already support structural refinement and
interface checks. Sealed existential witnesses add a generic
\emph{representation-independent} replacement guarantee. A private field alone
does not abstract the representation of a public value whose numeric, record,
or list structure remains observable.

\section{First-class packages and existential placement}
\example{Heterogeneous values with their own operations}{\Shigher{}+\Abs{}}
\par\noindent\begin{minipage}{\linewidth}
\begin{lstlisting}
#Showable: exists A {
    value: A
    show: func(A) -> string
}
p: seal #Showable with (A = int) {
    value: 7
    show: func(x: int) -> string: "\(x)"
}
q: seal #Showable with (A = string) {
    value: "hello"
    show: func(x: string) -> string: x
}
items: [p, q]
\end{lstlisting}
\end{minipage}\par
Each package has its own type witness and seal. A list of packages differs from
one package containing a homogeneous list.

\example{Eliminating a first-class package}{\Shigher{}+\Abs{}}
\par\noindent\begin{minipage}{\linewidth}
\begin{lstlisting}
describe: func(p: #Showable) -> string:
    (open p as (A, P) {out: P.show(P.value)}).out
text: map(describe, items)     // ["7", "hello"]
\end{lstlisting}
\end{minipage}\par
At a sealed interface, projections are mediated by the opened abstract view.
The checker does not independently instantiate the hidden type for \code{show}
and \code{value}.

\example{Sharing two values under one witness}{\Shigher{}+\Abs{}}
\par\noindent\begin{minipage}{\linewidth}
\begin{lstlisting}
#ComparablePair: exists A {
    left: A
    right: A
    equal: func(A, A) -> bool
}
compare: func(p: #ComparablePair) -> bool:
    (open p as (A, P) {out: P.equal(P.left, P.right)}).out
\end{lstlisting}
\end{minipage}\par
Separate existential packages for the two operands would require an explicit
sharing certificate before one comparator could consume both.

\example{An abstract codec interface}{\Sone{}+\Abs{}; checked effects}
\par\noindent\begin{minipage}{\linewidth}
\begin{lstlisting}
#CodecModule: exists Value {
    encode: func(Value) -> bytes
    decode: func(bytes) -> Value !decode
}
\end{lstlisting}
\end{minipage}\par
A round-trip law requires decoding an encoded value to return the corresponding
abstract value. Persisted wire compatibility additionally constrains the chosen
byte representation; hiding an in-memory representation does not silently
version an external protocol.

\Needspace{180pt}
\section{Mixed prefixes and generic modules}
The order of module binders expresses representation dependence:
\par\noindent\begin{minipage}{\linewidth}
\begin{lstlisting}
module(A): {
    exists State
    empty: State
    push: func(A, State) -> State
}
\end{lstlisting}
\end{minipage}\par
The logical order is $\forall A\,\exists State$. It allows the hidden semantic
state type to depend on $A$ while constraining one uniform module subject. A
standard witness is $State=\List(A)$, with empty list and a generic cons operation.
It differs from $\exists State\,\forall A$, which chooses one semantic state type
before the element type varies.

When different \emph{runtime} modules are required for different arguments, an
ordinary module-producing function supplies the module subject. Quantifiers
specify its interface; they do not themselves allocate a record at each type
assignment. This keeps generativity, runtime allocation, and logical dependence
separate and explicit.

\chapter{Checking, execution, and formal obligations}
\label{ch:checking}

\section{Typing judgments and operational subjects}
Use rigid type variables, term variables, and scoped existential witnesses in
$\Gamma$. Computation typing has synthesis and checking forms
\[
 \Gamma\vdash e\Rightarrow T,
 \qquad \Gamma\vdash e\Leftarrow T.
\]
Their conclusion constrains successful results. Value membership is separately
written $\Gamma\vdash_v v:T$. Subtyping certificates establish
$\Gamma\vdash_{\Static}T\le U$ and are sound for semantic inclusion. A compiler
may fail to establish such a certificate without proving the denotation empty.

Universal introduction checks the same erased computation under a fresh rigid
type variable. Universal elimination substitutes any well-sorted source type
permitted by the formation profile, including a polytype in \Shigher{}.
Existential elimination opens one rigid witness shared by all projections and
checks its escape condition. Scope levels record dependencies, not universes.

Bidirectional typing separates annotation checking from inference
\cite{bidirsurvey}. Explicit type abstraction and instantiation provide a route
to checking every annotated System F derivation. Inference of missing
impredicative arguments is an optional front-end service, never a condition on
the decidability of checking the supplied elaboration.

\section{Partial correctness of typed computation}
A lambda is a value even when its body may fail. The lambda rule checks the body
under its parameter types. Application checks its function and argument types
before using the result type. It never assumes success of the body.

\begin{theorem}[Partial-correctness soundness]\label{thm:soundness}
Assume sound primitive rules, sound subtyping certificates, scope-correct
elaboration, and the explicit-assertion rules of \cref{ch:residual}.
If $\Gamma\vdash e\Leftarrow T$, then in each environment satisfying $\Gamma$,
every successful complete result of the elaborated computation lies in
$\sem T$. If a resulting closure is certified at $A\Rightarrow B$, a successful
call with an admitted $A$-packet returns only a $B$-result. Failure and divergence
are permitted; a computation cannot produce a successful value by subsuming
failure to another type.
\end{theorem}
\begin{proof}
Induct on the derivation. Constants and structural values use membership;
subsumption uses inclusion. The lambda case establishes the conditional
universal in \cref{eq:arrow}, and application instantiates it. Meet intersects
successful completion predicates. Universal introduction varies the rigid
valuation without changing erased code; elimination uses
\cref{prop:subst}. Existential rules preserve their shared witness. Source
assertions forward only validated successes and otherwise fail or remain pending.
Strict evaluation premises and retained demand dependencies distinguish failed
computations from returned values.
\end{proof}
The theorem imposes no termination premise. A separately selected execution
profile may reject recursive code or permit a fixed-point operator with its
standard partial type. A derivation from an inconsistent data hypothesis is not
an implementation certificate: source checking uses verified variable types,
not arbitrary unresolved propositions as axioms for ex falso.

\section{Symbolic evaluation and delayed structure}
For a symbolic expression $e$, retain the successful-completion predicate
$\operatorname{CompObs}(e)$ and its demand and failure dependencies. The exact
normalization target is
\[
 \sem{N(e)}=\operatorname{CompObs}(e).
\]
Abstract shape and bound information may guide propagation while the exact
residual graph remains present. A candidate concrete result with a pending
condition is an observation under that condition, not an unconditional export.

CUE records and arguments can contain suspended computations. A required field
whose successful-completion predicate is proved empty makes the record's
completion predicate empty. The normalizer propagates that refutation through
its required enclosing fields; a suspended implementation must not indefinitely
hide an established contradiction. Optional impossible fields instead constrain
presence, according to ordinary CUE field rules.

Static expression checking nevertheless verifies every subexpression from its
own justified interface. A contradiction elsewhere in a record does not excuse
an invalid field selection or produce a certificate for an unrelated expression.
Within a function, a proof that the body fails under certain arguments is scoped
to those arguments. It does not turn the returned closure into a failed
closure-construction computation. The value/computation distinction for
non-strict semantic subtyping is relevant here \cite{nonstrict}.

The reference System F sublanguage uses strict application to value arguments.
A demand-driven evaluator can postpone computation, but commits a completed
call only when its required argument and result dependencies have been
validated. It must not use the empty-domain arrow identity to discard an
unvalidated argument and produce an arbitrarily typed result. Memoization
preserves these dependencies as well as code identity and captured bindings.

\section{Static certificates and executable assertions}
Every implementation-to-interface obligation has two destinations: a checked
certificate or a blocking diagnostic. This applies to structural, kind, literal,
polymorphic, callable-protocol, and dependent-result requirements alike.
An unresolved implication does not become an inserted runtime check.

An explicitly written assertion has a separately checked executable meaning.
Its successful-result predicate contributes to the body's inferred type, while
its execution can remain pending on data. The pending operation is source code,
not an unproved type obligation. This distinction permits cheap certification
of \code{x & >0} without proving that every possible \code{x} is positive.

An unspecified implementation may remain an ordinary data hypothesis after its
explicit interface passes relevant consistency and eager refutation. The
interface check does not wait for a body or an observing call. A client can be
checked under that interface;
linking verifies the actual implementation against that interface. A completed
executable artifact has neither unchecked imports nor outstanding annotation
obligations. The stages and minimum judgments are fixed in \cref{ch:residual}.

\chapter{Declarative checking and explicit assertions}
\label{ch:residual}
\section{The checking contract}
The checking profile \Static{} fixes a mandatory static fragment independently
of the formation profiles \Uone{}, \Sone{}, and \Shigher{}, and the extensions
\Dep{} and \Abs{}. Every annotated implementation receives a complete
certificate before executable use. Every explicit interface first passes the
relevant-consistency check of \cref{sec:relevant-consistency}, including an
interface with no implementation. Body refinements may follow from symbolic
reasoning or from the successful-result rules of assertions already in its body.

Elaboration produces
\[
 (e^\circ,\;\pi,\;R),
\]
where $e^\circ$ is the source program with its type annotations erased, $\pi$
certifies its interface and the source relevance check, and $R$ retains the exact data constraints and
execution dependencies of its explicit operations. No executable node in $R$
is created solely to satisfy an annotation. The observable outcomes of checking
are:
\begin{enumerate}
\item \emph{Accepted}: explicit interfaces pass relevant consistency and all
typing and interface implications have certificates; source assertions may still
await their data.
\item \emph{Blocked}: an invalid relevant observation, a missing annotation or
proof, an unsupported operation, an unproved implication, or exhausted static
resources.
\item \emph{Refuted}: a checked empty value or completion predicate, propagated
at its scope.
\end{enumerate}
A refuted required data subject is bottom. A refuted function-body path becomes
a failing computation at that path, rather than an unrelated global
contradiction. A blocking diagnostic does not assert semantic emptiness.

A module interface may remain open until linking. Refining that interface
triggers the same accumulated-conjunction check; it does not attach executable
checks to a cached closure. Completing missing declarations may resolve a
blocking diagnostic, but cannot remove a proved contradiction.

\section{The mandatory type fragment}
\label{sec:strict-fragment}
The strict grammar is finite, alias-expanded type syntax:
\[
\begin{aligned}
 T,U ::= {}& \bot\mid\top\mid\alpha\mid k\mid c
       \mid T\land U\mid T\lor U\mid T\Rightarrow U\\
       &\mid\{\ell_i^{m_i}:T_i\}_{i\in I}
       \mid[T_1,\ldots,T_n]\mid\List(T)
       \mid\forall\alpha\,T\mid\exists\alpha\,T.
\end{aligned}
\]
Here $k$ ranges over CUE kinds and $c$ over literals. Field mode $m_i$ records
requiredness or optionality, together with an explicit absence constraint when
present. The permitted quantifier positions are those of the selected formation
profile. Callable packets carry their statically declared arity, labels,
requiredness, defaults, and bound slots. Recursive aliases are outside this
finite kernel unless equipped with a separately specified contractive
regular-type procedure.

The fragment includes all of the following, not just record subtyping:
\begin{itemize}
\item kind membership and the kind lattice, including
$\Int\le\Num$, $\mathsf{Float}\le\Num$, and
$\Num=\Int\cup\mathsf{Float}$ and
$\Bool=\{\mathsf{false},\mathsf{true}\}$;
\item literal equality, literal membership in kinds, and literal inclusion in
finite unions and intersections;
\item record width, depth, presence, and required-field checks; fixed list
lengths and component types; homogeneous lists and their element types;
\item contravariant arrows, finite union/intersection reasoning, and scoped
universal and existential introduction and elimination.
\end{itemize}
The kind table distinguishes integers from decimal floating-point values,
strings from bytes, and data kinds from closures. Literal comparison and
arithmetic preserve CUE's specified numeric representation rules
\cite{cuespec}.

Consequently both
\[
 \{a:1,b:\mathsf{true}\}\le\{a:1\},
 \qquad (1\lor2)\le\Int
\]
are mandatory, although only the first is structural subtyping in the record
sense. A kernel cannot postpone either judgment to a runtime cast.

\subsection{Records are closed for typing}
The parameter's declared field inventory is closed for static reasoning:
unmentioned fields cannot be invented to justify selection or application.
Width subtyping remains available. The record type $\{a:1\}$ guarantees the
field $a$ and may describe a value also having $b$; only $a$ is accessible through
that static interface. This is distinct from \code{close({a: 1})}, which also
forbids extra fields in the data predicate.

Interpreting every bare type as an exact-label set would contradict the required
width relation, so the two notions of closure must remain separate. Extra fields
on an argument are permitted by width subtyping unless explicit data closedness
forbids them. Missing required fields are blocked, never guessed from an open
ambient schema.

\example{Closed static field inventory}{\Uone{}+\Static{}}
\begin{lstlisting}
get: func(r: {a: int}) -> int: r.a
ok: get({a: 1, b: true})        // 1: width subtyping

bad: func(r: {a: int}) -> int: r.b
// blocked at definition: b is absent from the parameter interface

needsB: func(r: {a: int, b: bool}) -> int: r.a
badCall: needsB({a: 1})         // blocked: required field b is missing
\end{lstlisting}
An explicit occurrence test can refine an optional field to present. A finite
union of record shapes is checked branchwise. A general dynamic-label map
requires its own declared map interface; it is not an excuse to residualize a
missing field in a finite record type.

\section{Strict typing rules}
The internal language records lambda annotations, type abstractions, type
applications, and the branch choices used by union/intersection derivations.
Source inference may insert them. The kernel checks the resulting proof graph
without searching for arbitrary types.

Representative mandatory rules are
\[
\frac{\Gamma,x:S\vdash e\Leftarrow T}
 {\Gamma\vdash\lambda(x:S).e\Leftarrow S\Rightarrow T}
\qquad
\frac{\Gamma\vdash f\Rightarrow S\Rightarrow T\quad
      \Gamma\vdash x\Leftarrow S}
 {\Gamma\vdash f(x)\Rightarrow T}.
\]
These are computation judgments: success has the displayed type. They do not
assert that evaluating the computation produces a value. The rules for explicit
polymorphism are
\[
\frac{\Gamma,\alpha:\Type\vdash e\Leftarrow T\quad
       \alpha\notin\FV(\Gamma,e_{\rm run})}
 {\Gamma\vdash e\Leftarrow\forall\alpha\,T}
\qquad
\frac{\Gamma\vdash e\Rightarrow\forall\alpha\,T\quad S:\Type}
 {\Gamma\vdash e[S]\Rightarrow T[S/\alpha]}.
\]
$S$ may contain quantifiers in \Shigher{}. Checking the substitution is mandatory;
inferring $S$ is not. The corresponding existential rules choose or open one
shared witness with its usual escape condition.

Intersection introduction checks the same expression at both types. Union
introduction records a proved injection; union elimination checks every
possible incoming branch. Subsumption requires a checked inclusion derivation.
A supplied literal synthesizes its singleton, and a supplied record synthesizes
its present fields. These precise descriptors are not silently widened before
checking a singleton or required-field obligation.

For unification expressions, the mandatory primitive rule is
\begin{equation}
\frac{\Gamma\vdash e\Rightarrow S\qquad
      \Gamma\vdash d\Rightarrow T}
 {\Gamma\vdash e\mathbin{\&}d\Rightarrow S\land T}.
\label{eq:meettyping}
\end{equation}
The result can fail. Both operands are checked; the possibility of a conflict
cannot conceal an ill-typed subexpression. A known local conflict is diagnosed
by the eager rules below. An explicit failure computation has type
$\Comp(\bot)$ and can occur in a function body or a deliberate failing branch.
It supplies no successful witness and permits no arbitrary cast of another
expression.

\example{Annotations assert properties of the body}{\Uone{}+\Static{}}
\begin{lstlisting}
id(A): func(x: A) -> A: x          // accepted
wrong(A): func(x: A) -> A: 0       // blocked: 0 is not arbitrary A
asText: func(x: int) -> string: x  // blocked: int does not imply string
asTwo: func(x: int) -> 2: 1        // refuted singleton obligation
stop(A): func(x: A) -> A: _|_      // explicit failure; no successful return
\end{lstlisting}
Every result annotation generates the same implication judgment. The body is
checked before result subsumption; failure of that implication blocks the
annotation. A source meet uses \cref{eq:meettyping}; no meet is added by checking.
The explicit failure term has no successful result and is checked by its own
primitive rule, after all enclosing structural obligations have been verified.

\subsection{A finite subtyping proof calculus}
Every kernel verifies reflexivity, transitivity with an explicit intermediate
type, bottom/top, the finite kind table, and literal membership. It verifies
record width/depth and presence rules, covariant products/lists, and
contravariant arrows. It also verifies the lattice rules
\[
 S\land T\le S,\quad S\le S\lor T,\quad
 \frac{R\le S\quad R\le T}{R\le S\land T},\quad
 \frac{S\le R\quad T\le R}{S\lor T\le R},
\]
finite distributivity, and the arrow laws in \cref{prop:arrow}. Quantified
inclusions carry explicit instantiations, openings, or the pointwise rules of
\cref{ch:refinement}. Bounded universal comparison uses the same bound under a
rigid variable, or a separately supplied proof of the required guarded
inclusion; unrestricted bounded-subtyping proof search is not a kernel rule.

This defines a proof-checkable extension of explicitly typed System F with
structural subtyping, unions, intersections, literals, and CUE kinds. Every
finite derivation whose explicit interfaces pass relevant consistency is accepted
when its proof graph fits the published resource bound. Inferred and intermediate
types may contain empty successful-result predicates; their use in a derivation
does not turn them into source interface declarations. It is not a completeness claim for every inclusion
between impredicative denotations. A smaller syntactic inclusion calculus can
be sound even when semantic inclusion has additional valid statements.

A practical front end memoizes type-pair obligations, normalizes finite record
fields and kind/literal unions, splits source unions and target intersections,
and tries the finitely recorded target-union/source-intersection alternatives.
Structural arrow checks recurse on smaller component types. Quantified
elimination uses explicit source instances or a finite candidate list; it never
searches an unbounded space of polytypes. Optional finite Boolean normalization
and the displayed distributive laws can derive more inclusions within the same
budget. A supplied proof can resolve a goal the search did not find. Explicitly
annotated higher-rank checking with union/intersection types and impredicative
applications has a closely related mechanized precedent \cite{jiang2025}.

\section{Required eager refutation}
\label{sec:eager}
Before issuing a certificate, run a shared eager-refutation service over the
available concrete and abstract descriptors. The service is also used by ordinary
CUE normalization. A conforming implementation publishes its finite transfer
rules and the conformance corpus for its chosen base-language version.
Type checking must preserve every refutation that this service produces; it
cannot downgrade one to a dynamic residual.

The mandatory minimum includes:
\begin{enumerate}
\item incompatible known kinds and unequal atomic literals, including numeric
kind distinctions and finite literal-union exclusion;
\item empty intersections of explicit constant numeric intervals, with open
endpoints and integer rounding handled correctly; concrete pattern checks;
\item known conflicting required fields, forbidden fields in explicitly closed
data, impossible field-presence combinations, and incompatible fixed list
lengths or components;
\item exact evaluation of enabled ground primitives and comparisons, invalid
known indexes, malformed call packets, duplicate labels, and missing arguments;
\item rejection of a finite alternative only after that alternative is refuted,
and rejection of their union only after all alternatives are refuted.
\end{enumerate}
Optional impossible fields describe absence and are not automatically an
inconsistent record. Required impossible fields propagate bottom. Branch and
quantifier assumptions remain attached to refutations; a counterexample under
an impossible guard proves nothing about a live branch. A contradiction in an
uncalled function body identifies a failing body path. It does not remove the
closure value or allow an unchecked subexpression to acquire an arbitrary type.

\example{Cheap contradictions and unresolved arithmetic}{Base fragment+\Static{}}
\begin{lstlisting}
eager: {x: <0, y: x & >0}    // refuted by constant-bound intersection

hypothesis: {
    a: int & b * 5
    b: a - 1
}                            // retained without symbolic elimination
concrete: hypothesis & {a: 0} // refuted: b = -1 and a = -5
\end{lstlisting}
The arithmetic relation has no integer solution: $a=5(a-1)$ implies $a=5/4$.
Its general emptiness need not be decided. Type checking establishes the kinds
of multiplication and subtraction, not the satisfiability of the equations.
Grounding a particular instance enables exact refutation. This matches the
distinction between incomplete cyclic constraints and concrete evaluation in
CUE \cite{cuecycles}.

The corpus requirement concerns the eager service and inputs available at the
same checking stage. An observation that needs future runtime data is checked
when those data arrive. An implementation cannot claim compliance by running
no transfer rules or by exhausting its budget before every trivial judgment.
Resource exhaustion blocks acceptance and is reported separately.

\section{Relevant consistency of explicit descriptions}
\label{sec:relevant-consistency}
An explicit interface describes successful observations on the input regions it
mentions. Arrow membership constrains partial computations; relevant consistency
checks whether that explicit description has already refuted a stated successful
observation. These are separate judgments. A relevance diagnostic blocks source
admission without equating the denotation of a function type to bottom.

\subsection{Relevant domains and pointwise result specifications}
For a finite conjunction of explicit arrow clauses
\[
 C=\bigwedge_{i\in I}(A_i\Rightarrow B_i),
\]
let $p$ range over complete call packets. Its \emph{relevant domain} is
$D_C=\bigcup_{i\in I}A_i$. Its partial result specification is
\[
 \mathcal B_C(p)=\bigcap_{\{i\mid p\in A_i\}}B_i
 \qquad(p\in D_C).
\]
Outside $D_C$ this specification is undefined: there is no relevant observation
there. Undefinedness is neither a top result nor a failed computation.

For each nonempty set $J\subseteq I$, define the exact input region and its
combined result predicate by
\begin{equation}
 R_J=\left(\bigcap_{j\in J}A_j\right)
       \setminus\left(\bigcup_{i\in I\setminus J}A_i\right),
 \qquad B_J=\bigcap_{j\in J}B_j.
 \label{eq:relevance-regions}
\end{equation}
The $R_J$ partition $D_C$; $\mathcal B_C$ is $B_J$ on $R_J$. With dependent
results, retain $B_j(p)$ in the same packet context. Calling conventions are
part of $A_i$, so arguments with different arities or incompatible required
labels need not inhabit overlapping regions.

Let $\Gamma\vdash_{\Static}P\equiv\bot$ denote a verified refutation by the
finite eager service and the selected strict proof rules. Failure to establish
this judgment is not a satisfiability proof. Complete the mandatory eager pass
before classifying a region; resource exhaustion blocks the check.

\begin{definition}[Relevant consistency]
\label{def:relevant-consistency}
An explicit arrow conjunction $C$ fails relevant consistency in declaration
context $\Gamma$ when a generated nonempty index set $J$ satisfies
\begin{equation}
 \Gamma\not\vdash_{\Static}R_J\equiv\bot,
 \qquad
 \Gamma,p:R_J\vdash_{\Static}B_J(p)\equiv\bot.
 \label{eq:relevant-invalid}
\end{equation}
A relevance certificate records the prescribed region checks and the absence
of such a refuted relevant observation. The predicate
$\operatorname{RC}_{\Static}(\Gamma,C)$ denotes passage of this finite check.
\end{definition}
For nondependent results the second judgment can be checked in $\Gamma$ alone.
A refuted input region is discarded before checking its result; inconsistency
of the region must not be used as an ex-falso proof of a live result conflict.
A domain that is merely unresolved is subject to the check.

This is a source well-formedness condition relative to a published refutation
kernel. Passage asserts neither totality nor satisfiability of every result
predicate. A body may still fail on particular inputs, or on all inputs. A
failure-only predicate synthesized for that body remains a valid internal type.

\example{Refuted explicit observations}{\Uone{}+\Static{}}
\begin{lstlisting}
f: (func(int) -> int) & (func(int) -> bool)
// blocked: on int, the explicit result is int & bool

g: func(int) -> _|_
// blocked: an explicit empty result on a non-refuted domain

h: (func(number) -> int) & (func(int) -> bool)
// blocked on int, even though the other numeric region is consistent
\end{lstlisting}
The first two declarations have the same relevant result specification. In the
third, every overlapping integer input requires an impossible successful result.
Successes elsewhere in the domain, or outside all declared domains, do not
satisfy that observation.

An explicit empty result is rejected by \cref{eq:relevant-invalid}; the internal
empty type, executable failure terms, and empty inferred results remain part of
the calculus. CUE's bottom expression retains its ordinary evaluation role
\cite{cuespec}.

\example{Nonconflicting overlap and disjoint cases}{\Uone{}+\Static{}}
\begin{lstlisting}
a: (func(number) -> number) & (func(int) -> int)
// accepted interface: integers on int, numbers on other numeric inputs

b: (func(int) -> string) & (func(string) -> int)
// accepted interface: the two domains are disjoint

c: (func(int) -> (0 | 1)) &
   (func(int) -> (1 | 2)) &
   (func(int) -> (0 | 2))
// blocked: all three results have empty intersection on int
\end{lstlisting}
The third example requires joint conjunction, not merely pairwise consistency:
each pair of its result predicates intersects. Equal-domain grouping exposes
the contradiction without enumerating input values.

\subsection{Source boundaries, rigid binders, and instantiation}
A \emph{relevance root} is the accumulated explicit type description of one
subject at one declaration boundary. Inline and independent annotations share
that root. It carries source origins, quantifier dependencies, and guards.
Parametric aliases expand before checking and retain the origin of their use.
An inferred type, a proof intermediate, or a selected generic instance creates
no additional root.

At a universal declaration, check its body under fresh rigid type variables and
the declared bounds. Refutation of a result must follow under those variables;
the checker does not choose a type assignment solely to make the result empty.
This convention permits declarations whose individual instances may fail.

\example{A generic computation and an explicit specialization}{\Uone{}+\Static{}}
\begin{lstlisting}
meet(A, B): func(x: A, y: B) -> (A & B): x & y
// accepted: A & B is not refuted at the declaration

good: meet({a: 1}, {b: true}) // {a: 1, b: true}
bad:  meet[int, bool](1, true) // _|_, a failing computation

specialized: func(int, bool) -> (int & bool)
// blocked: a newly declared interface with refuted results
\end{lstlisting}
The instance $[\Int,\Bool]\Rightarrow\bot$ in the call is derived from an
accepted generic declaration. It is an internal description of the computation,
not a new assertion that this particular interface describes successful behavior.
Instantiation and inference must not convert it into an explicit annotation and
rerun relevance. An annotation written by the programmer does create a root,
including one written through a parametric alias.

Conjoining a universal scheme with another explicit arrow also requires checking
their stated overlapping obligations. The checker uses the non-generic domains
of sibling clauses as finite \emph{anchors}. Match a universal clause's domain
pattern against an anchor, verify the resulting substitution and bounds, and
add that instance to the root's observation checks. Only variables bound by
that universal may be substituted; outer existential dependencies remain
correlated variables. The minimal rule uses finite structural matching, with
repeated variables required to match the same type. Additional certified matches
may be supplied within the same budget.

\example{A universal obligation at an explicit sibling domain}{\Uone{}+\Static{}}
\begin{lstlisting}
quiet(A): func(A) -> A
quiet:    func(int) -> bool
// blocked before any body is supplied
\end{lstlisting}
The integer domain anchors the instance $A=\Int$, giving the conflicting
results $\Int$ and $\Bool$. Matching is triggered by the independently stated
integer interface, not by a call. The generic \code{meet} declaration has no
such independently stated ground-domain anchor. Its later applications do not
modify its source root.

Alpha-renaming, quantifier floating, and copied declarations retain which
variables a source clause depends on. An integer clause floated under
\code{forall A} remains independent of $A$ and remains an anchor. Equivalent
same-domain result splitting and conjunction regrouping generate identical
relevance checks. Arbitrary substitution of an accepted generic declaration by
a newly written specialized interface need not preserve source admissibility;
its denotation still obeys the substitution lemma.

\subsection{Nested observations, alternatives, and guards}
Traverse explicit parameter, result, record, and list interfaces recursively.
At each function-valued subject, collect conjunctive arrow requirements before
forming regions. Required-field and element-presence guards are inherited from
the enclosing description. An impossible optional field constrains absence;
an empty-list alternative makes no element observation.

Union alternatives describe different possible interfaces. Check each surviving
explicit alternative separately; never intersect arrow results from opposite
sides of a union. A guard proved impossible by ordinary eager normalization
removes its branch before relevance traversal. A semantically inhabited but
relevance-invalid arrow alternative cannot be silently discarded as though it
were an empty data alternative.

\example{Separate alternatives and nested callback contracts}{\Uone{}+\Static{}}
\begin{lstlisting}
choice: (func(int) -> int) | (func(int) -> bool)
// accepted interface: alternatives, not simultaneous obligations

callback: func(
    (func(int) -> int) & (func(int) -> bool),
) -> string
// blocked: the explicit callback interface is relevance-invalid
\end{lstlisting}
The source checker validates stated observations even in negative type
positions. This is a formation condition, independent of the semantic variance
of that position.

For dependent results use the packet context in \cref{eq:relevant-invalid}.
For example, an integer parameter $x$ with result constraints $y>x$ and $y<x$
has refuted successful results whenever the selected arithmetic rules establish
the contradiction. Unresolved arithmetic remains unresolved; relevance does not
supply a general satisfiability oracle.

\subsection{Finite checking and preservation properties}
A straightforward algorithm enumerates the regions of
\cref{eq:relevance-regions}, discards refuted guards, and intersects the remaining
result predicates. A shared decision DAG avoids duplicating equivalent regions.
Equal-domain aggregation, disjoint-kind tests, finite literal sets, and constant
interval intersections are mandatory fast cases. Nonempty subsets of three or
more clauses must be considered when their cumulative results can conflict.

A conflicting subset can guide the search. For $J\subseteq I$, put
$Q_J=\bigcap_{j\in J}A_j$. If its result conjunction is refuted, every
surviving refinement of that overlap has the same conflicting requirements.
Split $Q_J$ by the remaining domain guards and stop at a non-refuted exact
region. A certificate retains that region, its contributing source origins, and
the result refutation. A known inhabitation witness for $Q_J$ can identify a
live region directly. This search need not enumerate unrelated regions;
unknown overlap alone is not a proof that any particular exact region survives.

All expansion, matching, guard checks, and region splitting consume the static
budget. Exhaustion blocks admission; it never omits a required relevance check
or turns it into a runtime assertion. On a finite closed kind/literal partition,
the reference procedure decides relevant consistency exactly. Richer predicates
use the published finite refutation relation. The existing rank profiles and
\Static{} remain orthogonal: every profile applies relevance at explicit
interface boundaries, while their available match and proof forms differ.

\begin{proposition}[Relevance coherence]
\label{prop:relevance-coherence}
For fixed source roots, scopes, and canonical checking rules, permuting or
duplicating interface conjuncts preserves relevant consistency. Replacing a
same-domain arrow with the conjunction of its result-split arrows preserves
the relevant regions and their result predicates. These transformations do not
change erased code.
\end{proposition}
\begin{proof}
The relevant domain is a union of declared domains; each relevant result is
an intersection of all applicable requirements. Both operations are
associative, commutative, and idempotent. Result splitting repeats the same
guard, and alpha-renaming preserves its dependencies. Canonical sharing makes
the prescribed comparisons and their budget charges independent of grouping.
The relevance pass constructs evidence or diagnostics, never executable nodes.
\end{proof}

Relevant consistency is a source admission condition, not a predicate equality
in $\Pow(V)$ and not an additional lattice operation. Adding a declaration can
expose an invalid relevant observation while leaving the denotation inhabited.
If a conflicting region has an inhabitation witness, further arrow conjuncts
cannot remove the conflict: they only refine results there. When the region is
merely non-refuted, the blocking judgment is relative to its retained guard and
checking context; it is not a proof of semantic inconsistency. The normalizer
therefore never turns a relevance diagnostic alone into bottom. Denotational
refinement and annotation erasure retain their existing theorems.

\section{The static--dynamic boundary}
\label{sec:dynamic-admission}
Annotations and executable assertions have different judgments. An annotation
$T$ requires a derivation of the body's successful-result inclusion in $T$.
The expression \code{e & P} executes a constraint operation written by the
programmer. When $P$ has a checked first-order interpretation, use
\begin{equation}
 \frac{\Gamma\vdash e\Rightarrow S\qquad
       \Gamma\vdash P\ \mathsf{predicate}(K)\qquad
       \Gamma\vdash S\le K}
      {\Gamma\vdash e\mathbin{\&}P\Rightarrow
           \{v\in S\mid P(v)\}}.
 \label{eq:asserttyping}
\end{equation}
Here $K$ is the predicate operation's admissible data interface. A predicate
with ordinary CUE meet semantics on all kinds, such as a kind constraint, can
use $K=\top$ and fail on incompatible kinds. Operation checking is static;
success or failure of the explicit assertion may depend on runtime data.

The predicate rule and the ordinary meet rule are source primitives, not casts
introduced by a subsumption rule. For \code{x: int}, the body \code{x & >0}
synthesizes $\Int\cap\sem{>0}$. Checking that result against the same bound is
reflexivity. Checking the body \code{x} against that bound instead requires
$\Int\subseteq\sem{>0}$ and is blocked.

The minimum entailment kernel includes conjunction elimination/introduction,
reflexivity of alpha-equivalent dependent predicates with the same captured
dependencies, constant-bound inclusion, and the strict rules above. These rules
accept explicitly asserted refinements without a general arithmetic theorem
prover. More involved implication proofs may be supplied or searched within a
budget. An unproved implication always blocks; a runtime assertion must already
occur in the source term.

A source predicate can retain a pending execution exactly when:
\begin{enumerate}
\item its operands and data interface have been checked;
\item its successful-result rule is certified independently of the requested
annotation;
\item its executable predicate and all shared dependencies remain in the
source evaluation graph;
\item the eager service has not already proved its completion predicate empty
in the applicable scope.
\end{enumerate}
A predicate returns pass, fail, or pending. Only pass releases a completed value
with its refinement. A refuted predicate is normalized to failure instead of
remaining pending. A predicate not refuted by the finite abstraction is a
possible residual, not a proved satisfiable relation.

Bounds, regular expressions, finite data traversals, and comparisons on actual
arguments admit such rules. Quantified function behavior, impredicative
membership, and opaque representation membership require static evidence;
their mere occurrence in an annotation supplies no executable test.

\example{An annotation and an explicit assertion}{\Uone{}+\Dep{}+\Static{}}
\begin{lstlisting}
positive: func(x: int) -> (int & >0): x
// blocked: int does not imply >0

positiveChecked: func(x: int) -> (int & >0): x & >0
// accepted: the body explicitly establishes the refinement on success
ok:  positiveChecked(3)   // 3
bad: positiveChecked(0)   // _|_
\end{lstlisting}
The accepted body has the same assertion before and after certification. Its
annotation erases; its \code{& >0} operation remains.

\example{Explicit kind refinement}{\Uone{}+\Static{}}
\begin{lstlisting}
pick: func(x: int | string) -> int: x
// blocked: the result inclusion is false

pickChecked: func(x: int | string) -> int: x & int
// accepted: successful unification produces an integer

same(A: number): func(x: A) -> A: x + 1
// blocked: a numeric upper bound does not establish closure of A
\end{lstlisting}
The source assertion in \code{pickChecked} can reject a string. A type
annotation on \code{pick} cannot introduce that rejection.

\section{Conjunctive declarations and body evidence}
\label{sec:declaration-coherence}
For each subject, collect a scoped conjunction of its interface predicates and
its implementation clauses. Inline and separate function-type annotations enter
the same conjunction. Their placement, order, and parenthesization have no
operational significance. A concrete implementation contributes its original
body, captured environment, and call protocol.

Let $\mathcal C_f$ be the finite set of interface conjuncts for $f$, and let
$e_f$ be its body. First check relevant consistency of the accumulated explicit
root, irrespective of whether $e_f$ is available. The checker then verifies $e_f$ against
$\bigwedge\mathcal C_f$, checking each universal clause at its rigid instance.
It retains the body's inferred refinement or replayable derivation, rather than
using a weaker declared result type as its only evidence. A separately compiled
opaque import instead exposes precisely its certified interface. A stronger
client requirement needs stronger evidence from that module, not a client-side
wrapper.

\example{Five presentations of one implementation constraint}{\Uone{}+\Dep{}+\Static{}}
Each listing is an independent declaration set for the same field.
\begin{lstlisting}
positive: func(x: int) -> (int & >0): x & >0
\end{lstlisting}
\begin{lstlisting}
positive: func(int) -> (int & >0)
positive: func(x: int) -> _: x & >0
\end{lstlisting}
\begin{lstlisting}
positive: func(x: int) -> int: x & >0
positive: func(int) -> >0
\end{lstlisting}
\begin{lstlisting}
positive: func(int) -> int
positive: func(x: int) -> >0: x & >0
\end{lstlisting}
\begin{lstlisting}
positive: func(int) -> int
positive: func(int) -> >0
positive: func(x: int) -> _: x & >0
\end{lstlisting}
All five have the erased body $\lambda x.\,(x\mathbin{\&}\sem{>0})$ and the
combined obligation
\[
 (\Int\Rightarrow\Int)\cap(\Int\Rightarrow\sem{>0})
 =\Int\Rightarrow(\Int\cap\sem{>0}).
\]
The body synthesis and result inclusion are identical, modulo redundant
conjuncts. The relevant integer region has result $\Int\cap\sem{>0}$, which is not
refuted. All five roots pass relevant consistency. Each declaration set is
accepted; calls at $3$ return $3$ and calls at $0$ fail.

\begin{proposition}[Declaration coherence]
\label{prop:decl-coherence}
For a fixed implementation origin and scope graph, regrouping, permuting, or
duplicating interface conjuncts preserves the accumulated predicate and erased
body. Splitting a result intersection into separate same-domain arrow
constraints preserves both and preserves relevant consistency by
\cref{prop:relevance-coherence}. The required kernel rules certify the same result
for the five presentations above. Separately written examples are compared
up to a consistent renaming of their implementation origins.
\end{proposition}
\begin{proof}
Predicate conjunction is associative, commutative, and idempotent. The arrow
intersection law in \cref{prop:arrow} distributes a fixed domain over result
intersection. Body extraction ignores result annotations, and checks all
conjuncts against the same retained body derivation. No checking step contributes
an executable operation.
\end{proof}
For arbitrary equivalent formulas, bounded proof search need not discover the
same derivation. The mandatory syntactic regroupings use canonical conjunction
DAGs and shared body evidence, so they neither duplicate the budget charge nor
depend on file traversal order.

\section{Failure-only predicates and their scope}
\label{sec:failure-propagation}
Write $\operatorname{Succ}_\Gamma(e)$ for the successful-completion predicate
of $e$ in its retained environment. The finite kernel recognizes empty kind
intersections, incompatible constants and bounds, and the empty result of an
applicable arrow-clause intersection. A proved empty predicate triggers
\[
 \Gamma\vdash\operatorname{Succ}(e)\subseteq\varnothing
 \quad\Longrightarrow\quad e\longrightarrow\code{_|_}.
\]
This is mandatory whenever the implemented eager rules establish the premise.
No supplied body or concrete value is needed to refute an already empty data
hypothesis. Required field propagation turns a bottom field into an inconsistent
configuration. A lambda body is a different scope: its failure describes a call,
not a failure to construct the closure.

For an internal result specification
$(A\Rightarrow B)\cap(A\Rightarrow C)$, any correctly admitted call at an
argument in $A$ has successful results in $B\cap C$. When that intersection is
empty, the call reduces to bottom. This rule applies to derived generic
instances and inferred body paths. A newly written interface asserting the same
empty relevant result is instead blocked by \cref{sec:relevant-consistency},
before a call or implementation is required.

\example{Invalid interfaces and scoped computation failure}{\Uone{}+\Static{}}
\begin{lstlisting}
q: func(int) -> int
q: func(int) -> bool        // blocked: explicit relevant conflict

meet(A, B): func(x: A, y: B) -> (A & B): x & y
observed: meet(1, true)     // _|_: computation, not interface rejection

stop(A): func(x: A) -> A: _|_
// a closure with a failing body and a non-refuted generic interface
\end{lstlisting}
Type checking still verifies subexpressions and calling protocols before using
an empty-result rule. An ill-typed operation is not repaired by assuming that
some unrelated computation fails.

\section{Stages and mandatory certificates}
\begin{longtable}{@{}P{0.18\textwidth}P{0.75\textwidth}@{}}
\toprule
Stage & Obligations completed before advancing\\\midrule
\endfirsthead
\toprule Stage & Obligations completed before advancing\\\midrule
\endhead
Elaboration & Scope and sort correctness; profile restrictions; alias expansion;
fixed calling protocols; accumulation of interface conjuncts and identification
of explicit source assertions.\\
Interface admission & Relevant consistency of every explicit root, including
unimplemented imports and nested callback types; guarded region checks and
source-anchored universal matching.\\
Definition checking & Every function body, including uncalled bodies, has a strict
derivation. Prove every result implication using the unchanged body, including
the successful-result rules of its explicit assertions.\\
Linking & Every supplied user function, builtin, foreign adapter, and captured
callable satisfies the required strict interface. Unspecified imports remain
link obligations and prevent a closed executable artifact.\\
Call checking & Check every statically present application and type application.
Prove packet compatibility and every argument/result implication. Source
assertions in arguments may supply the required refinements. Generic instances
are derived types, not new relevance roots. Newly visible literal and shape
contradictions are propagated at their scopes.\\
Normalization & Run eager transfers and exact enabled primitives; maintain joint
constraints, alternatives, guards, and source-assertion dependencies. Propagate
established empty predicates at their scopes.\\
Concrete observation & Demand and discharge every source assertion on the
observed path and every required data-validation dependency. Failure produces bottom;
pending conditions prevent unconditional completed export.\\
\bottomrule
\end{longtable}
Modules may store relevance-checked unlinked declarations as hypotheses. Their certificates are
conditional on named imports, as a typing derivation is conditional on its
context. A body or import does not become certified merely because no call
currently observes it. Calling an unknown function does not synthesize a suitable
implementation from the desired codomain.

\example{Higher-order obligations are never runtime residuals}{\Uone{}+\Static{}}
\begin{lstlisting}
apply: func(g: func(number) -> string) -> string: g(1.5)
ints: func(x: int) -> string: "ok"
blocked: apply(ints)         // blocked before execution

source: func(int) -> string // an unresolved external implementation
client: func() -> string: source(3)
\end{lstlisting}
The strict contravariant rule rejects \code{ints} as the required callback.
The semantic possibility that its out-of-domain calls fail does not discharge
that rule. This is the same intentional separation between denotational safety
and a useful syntactic type discipline that prevents arbitrary dead code from
being accepted solely on a conjecture that it never returns.

\section{A bounded implementation}
Store types, interface conjunctions, and body summaries as hash-consed DAGs;
identify binders and implementations by stable origins. Store proof DAGs
separately from source evaluation nodes. The checker for one accumulated subject
has the following form:
\begin{lstlisting}[language={}]
check_subject(declarations, budget):
    scoped = canonicalize_scopes_and_conjuncts(declarations)
    body, interface, hypotheses = separate_code_and_predicates(scoped)
    state = eager_normalize_with_scoped_failure(
        conjunction(hypotheses, interface), budget)
    if state.required_subject_is_empty:
        return Refuted(state.proof)
    relevance = check_explicit_relevance(scoped.roots, state, budget)
    if relevance is invalid or relevance is exhausted:
        return Blocked(relevance.diagnostic)
    if body is absent:
        return Hypothesis(state, interface, relevance)
    source = elaborate_source_operations(body)
    evidence = check_body_and_explicit_assertions(source, budget)
    for clause in canonical_clauses(interface):
        proof = check_implication(evidence, clause, budget)
        if proof is absent:
            return Blocked(clause, reason="annotation not established")
        verify_proof(proof, budget)
    return Accepted(erase_annotations(source), evidence, relevance, state)
\end{lstlisting}
Elaborating a source assertion lowers an operation already present in the term;
it never consults the requested output predicate to invent a check. A missing
body remains an ordinary hypothesis, not a certified implementation. Linking a
body reruns the same interface judgment. A retained source operation may await
its operands only after its operation rule and annotation implications have
been checked.

All work is charged, including alias expansion, substitution, arithmetic, and
Boolean splitting. Canonicalized duplicates are shared. Published budgets and
deterministic traversal make required regroupings reproducible. Budget
exhaustion blocks certification; it does not alter code or create a runtime
obligation. Source evaluation has its own resource policy and cannot report a
false validation pass.

Let $L$ be the size of a supplied proof DAG after sharing and let $C$ bound its
charged local comparisons. Verification costs $O(L\,C)$ time and $O(L)$ graph
space. Search uses a declared fuel $B$, costing at most $O(B\,C)$ work.
Unions and intersections can require exponentially many alternatives. The
kernel promises its finite rules and minimum fixtures, not complete semantic
inclusion at a fixed budget.

The minimum direct rules cover kinds, constants, finite record lookup, tuple
projection, explicit type substitution, meet, predicate reflexivity, and the
five declaration presentations. In particular, explicit \code{x & >0} needs
no proof that all inputs satisfy \code{>0}. Its primitive successful-result rule
supplies the fact needed by annotation checking. Explicit integer-to-integer
and integer-to-boolean conjuncts are rejected by relevance before any body check.
A pass-through implementation that stores such an interface as a hypothesis is
not conforming.

\section{Erasure and refinement preservation}
\label{sec:erasure}
Let $|e|$ erase function-type annotations, type abstractions, type-instantiation
selections, and proof arguments. It preserves the source's data constraints,
term-level meets, predicate operations, captures, labels, omission behavior,
and computational indices. Representative clauses are
\[
\begin{aligned}
 |\operatorname{func}(x:S)\mathbin{\to}T:e|
   &=\lambda_{\mathit{protocol}}x.\,|e|,\\
 |\forall A\,e|&=|e|,& |e[T]|&=|e|,\\
 |e\mathbin{\&}P|&=|e|\mathbin{\&}|P|
     &&\text{for an executable assertion.}
\end{aligned}
\]
The same predicate spelling can occur in an annotation or as an operand of a
source operation. Its elaborated role determines erasure: \code{-> >0} is a
proof requirement; \code{x & >0} computes a constrained result. Ordinary CUE
field constraints remain part of the data language, including their validation
dependencies. They are not discarded as function-type annotations.

\begin{theorem}[Annotation erasure]
\label{thm:erasure}
Assume the source-operation rules preserve their specified observations and
the static kernel is sound. For an Accepted implementation in an environment
satisfying its certified interface, erasing its annotations preserves successful
returns, explicit assertion failures, and the retained demand dependencies.
The erasure is independent of the placement and grouping of its annotations.
\end{theorem}
\begin{proof}
Each typing rule contributes evidence and no executable filter. Lambda erasure
retains the binder's operational protocol and the erased body. Instantiation
changes only the proof environment. Subsumption and intersection introduction
retain the same term. Relevant-consistency checking contributes only static
evidence or a blocking diagnostic. Executable meets and assertions have the same evaluation
rule on each side of erasure. Induction on evaluation, together with
\cref{prop:decl-coherence}, gives observation preservation and declaration
independence. Computational witnesses, sealing operations, and omission defaults
remain terms and follow their specified source rules.
\end{proof}

For scoped declarations $D$ and $E$ sharing an interface, the elaborated
predicate is refinement-homomorphic:
\begin{equation}
 \sem{D\mathbin{\&}E}=\sem D\cap\sem E.
 \label{eq:declaration-homomorphism}
\end{equation}
For a common body origin, the executable contribution on both sides is the
same erased body; interface conjunction adds predicates and proof obligations.
An incompatible additional annotation refutes a predicate, fails relevant
consistency, or blocks its conformance certificate. It never edits execution to enforce the new annotation. Distinct
body origins retain the closure-identity semantics in \cref{ch:programs}.
Equation \eqref{eq:declaration-homomorphism} concerns declaration predicates;
application still satisfies the inclusion \eqref{eq:appmeet}, not a general
meet-homomorphism equation.

Keep the exact residual predicate and its scope graph during normalization:
\begin{equation}
 P\longrightarrow Q\quad\Longrightarrow\quad
 P\simeq Q,\qquad P\land R\simeq Q\land R.
 \label{eq:residual-equivalence}
\end{equation}
A source assertion passing on a value establishes its refinement. Unresolved
captured dependencies remain attached to that success; for example,
\code{3 & (x + 1)} cannot become unconditional \code{3} while \code{x} is
unresolved. A proved empty required predicate propagates bottom. Quantifier
and branch scopes prevent a local body failure from becoming a false global
refutation.

\begin{theorem}[Static gate and explicit-assertion partial correctness]
\label{thm:residual-sound}
Assume sound strict rules, sound scoped eager refutation, dependency-preserving
normalization, and the source assertion rules. Every Accepted artifact has a
certificate for every annotation and has passed the required relevance checks
at each explicit interface root. In an environment satisfying its imports,
every successful completed observation has its stated type and demanded
refinements. Further conjunction cannot repair a proved data contradiction or
replace an atomic completed observation by a different atom. Erasure preserves
the implementation's observations.
\end{theorem}
\begin{proof}
The admission gate establishes the relevance-check clause; it contributes no
additional execution rule. Induction on body and implication certificates gives
\cref{thm:soundness}. The source assertion rule constrains successful results without a proof of
success. Equation \eqref{eq:residual-equivalence} preserves exact constraints,
and \cref{thm:monotone} lifts this preservation through quantification. A subset
of an empty predicate is empty; a subset of a singleton either retains that
atom or is empty. Demand dependencies prohibit unconditional release of a
conditional value. \Cref{thm:erasure} supplies the final assertion.
\end{proof}

\section{Profile combinations and conformance examples}
\begin{tabularx}{\textwidth}{@{}lY@{}}
\toprule
Combination & Mandatory checking organization\\\midrule
\Uone{}+\Static{} & Prenex schemes, monomorphic argument interfaces, finite
instance inference; explicit instances where local search is insufficient.\\
\Sone{}+\Static{} & The preceding kernel plus module-boundary existential
witnesses, sealing, and shared rigid openings.\\
\Shigher{}+\Static{} & Annotated arbitrary-rank types and explicit impredicative
instantiation. Quantified membership never becomes a runtime predicate.\\
\Dep{}+\Static{} & Mandatory strict skeleton; eager known dependent
contradictions; source assertions supply certified refinements, and every
annotation implication is proved.\\
\Abs{}+\Static{} & Strict abstract operations and scope checks. Representation
independence additionally uses an outcome-preserving simulation.\\
\bottomrule
\end{tabularx}

The following are minimum conformance obligations, not optional examples:
\begingroup
\small
\begin{longtable}{@{}P{0.61\textwidth}P{0.32\textwidth}@{}}
\toprule
Judgment or program & Required outcome\\\midrule
\endfirsthead
\toprule Judgment or program & Required outcome\\\midrule
\endhead
$(1\lor2)\le\Int$, $\Int\le\Num$ & Proved statically.\\
$(1\lor\mathsf{true})\le\Int$ & Blocked; no inserted kind check.\\
$\{a:1,b:\mathsf{true}\}\le\{a:1\}$ & Proved statically.\\
$\{a:1\}\le\{a:1,b:\Bool\}$ & Blocked required-field judgment.\\
Identity at $\forall A.\,A\Rightarrow A$ & Accepted with a rigid-variable proof.\\
Identity instantiated at its own quantified type & Accepted in \Shigher{} with
explicit instantiation.\\
\code{func(xs: [A,A]) -> A: xs[0] & xs[1]} & Accepted by tuple projection and meet.\\
An integer result annotated as string & Blocked at definition.\\
An integer callback used where a number callback is required & Blocked at the
static call or link boundary.\\
An explicit ground interval contradiction & Refuted before execution.\\
An unresolved arithmetic cycle with compatible operand kinds & Retained as data
constraints; ground instances checked.\\
A dependent result bound absent from the body evidence & Blocked unless its
implication is proved.\\
\code{x & >0} at \code{int & >0}, under \code{x: int} & Accepted by the explicit
assertion rule.\\
The five presentations in \cref{sec:declaration-coherence} & Accepted with one
erased body and equivalent contracts.\\
Explicit $\Int\Rightarrow\Int$ and $\Int\Rightarrow\Bool$ on one subject &
Blocked by relevance without a body or a call.\\
An explicit $\Int\Rightarrow\bot$ & Blocked by relevance.\\
A conflicting overlap with a consistent region elsewhere & Blocked; other
regions do not repair the overlap.\\
The generic meet declaration with rigid $A,B$ & Accepted; $A\cap B$ is not
refuted at the declaration.\\
A failing derived instance of generic meet & Reduced to bottom at its call
scope; no new explicit-interface relevance check.\\
An admitted derived call with a proved empty result intersection & Reduced to
bottom at the call scope, before implementation linkage if possible.\\
An unresolved polymorphic conformance judgment & Blocked, never a dynamic residual.\\
Static resource limit exceeded & Blocked with a resource diagnostic.\\
\bottomrule
\end{longtable}
\endgroup

\chapter{Feature correspondence and related designs}
\label{ch:comparison}
\section{The comparison baseline}
The July 2026 function proposal describes opaque function values, packet binding,
defaults, open signatures, and partial applications. The September theory
addendum analyzes a covariantly tightened call-constraint interpretation and
pointwise conjunction of bodies. It separately proposes coverage obligations
and compatibility checking \cite{cueproposal,cueaddendum}.

The present calculus expresses call restrictions existentially and reusable
capabilities universally. Its set lattice therefore includes both relations
on actual call cells and predicates on implementations. Their different
variance follows from their quantifiers.

\section{Feature coverage}
\small
\begin{longtable}{@{}P{.30\textwidth}P{.64\textwidth}@{}}
\toprule
Feature & Interpretation or construction in quantified CUE\\\midrule
\endfirsthead
\toprule Feature & Interpretation or construction in quantified CUE\\\midrule
\endhead
Legacy records, lists, unions, bounds & Unchanged data predicates and legacy elaboration; \cref{thm:conservative}.\\
Definitions, closedness, optional and required fields & Retained field-presence and label constraints; quantifiers do not add fields.\\
References, embedding, comprehensions & Existing copying plus capture-avoiding binder and witness rebinding.\\
Literal bodies and inferred result constraints & Operational closures with a checked declared or inferred result predicate.\\
Positional and labeled arguments & Complete call packets and a fixed binder.\\
Required, optional, hidden, and anonymous parameters & Presence alternatives, package-qualified labels, and positional slots.\\
Body-only aliases & Lexical aliases with no public packet label.\\
Parameter order and uniqueness & Protocol well-formedness before value unification.\\
Parameter attributes & Retained metadata, interpreted only by the relevant declared extension.\\
Omission defaults & Implementation-owned packet completion; supplied arguments bypass them.\\
Caller disjunction defaults & Existing preference rules on the supplied value.\\
Default contributed by another declaration & Explicit defaulted adapter or a refinable description of a yet-unselected implementation protocol.\\
Label contributed to an existing closure & Explicit label adapter; fresh builtins may receive their generated binder at construction.\\
Function types and intersections & Sets of operational inhabitants satisfying all universal arrow clauses.\\
Open signatures & Existential protocol rows retained until the remaining required packet shape is known.\\
Higher-order parameters and results & Ordinary arrows; universals specify reuse across types.\\
Partial application & A new closure with normalized original-slot bindings and a residual contract.\\
Constraint on a partial value & An ordinary constraint on that residual closure.\\
Distinct bodies combined pointwise & An explicit body that calls both and unifies their result descriptions.\\
Self-unification and captures & Stable origins and structurally constrained captured environments.\\
Validators & Predicates on an existential subject, definable with existing structs.\\
List-based variadic use & One list-valued packet component, with a generic element type.\\
Nonrecursive execution and cycles & The legacy execution profile; structural recursion is a separate opt-in policy.\\
Builtins and foreign calls & Typed descriptors with an explicit admitted effect set and the existing bridge.\\
Compatibility checking & Packet-domain inclusion, output/effect refinement, and module simulations where abstraction applies.\\
Export, formatting, and round trips & Source export retains binder scopes, closure environments, and bound slots; data export retains its existing supported values.\\
Memoization and dependency tooling & Scope-preserving graph traversal and dependency-aware caching.\\
\bottomrule
\end{longtable}
\normalsize

Experimental operations that transform a closure elaborate to explicit
transformations. Contract intersection asserts properties of the original
closure.

\section{Three explicit translations}
\example{Supplying a label}{\Uone{}}
\par\noindent\begin{minipage}{\linewidth}
\begin{lstlisting}
raw: func(_~n: int) -> int: n + 1
named: func(x: int) -> int: raw(x)
a: named(x: 2)              // 3
\end{lstlisting}
\end{minipage}\par
The adapter's packet language admits \code{x}. Raw positional behavior remains
unchanged.

\example{Restricting public calls across files}{\Uone{}}
\par\noindent\begin{minipage}{\linewidth}
\begin{lstlisting}
raw: func(x: number) -> string: "\(x)"
call: {arg: number, result: raw(arg)}
call: {arg: int, result: string}

ok:  (call & {arg: 7}).result   // "7"
bad: (call & {arg: 1.5}).result // _|_
\end{lstlisting}
\end{minipage}\par
This is ordinary CUE instantiation and projection. The argument restriction is
an existential predicate, while \code{raw}'s capability remains universal.

\example{A first-order dependent result check}{\Uone{}+\Dep{}+\Static{}}
\begin{lstlisting}
legacyBody: func(int) -> int
checked: func(x: int) -> (int & >0): legacyBody(x) & >0
\end{lstlisting}
The body explicitly constrains the integer returned by \code{legacyBody}.
Its meet rule proves the refined codomain without assuming success of the
comparison. The type annotation erases and the source assertion remains.
Linking separately verifies \code{legacyBody}'s integer interface. Omitting the
source assertion requires an independent proof of its positive-result property.

\begin{proposition}[Contract intersection]
Under the universal arrow interpretation, conjoining a contract with a concrete
closure cannot change its call protocol or result. It preserves the singleton
closure or makes its intersection empty.
\end{proposition}
\begin{proof}
A subset of the singleton $\{f\}$ is either $\{f\}$ or empty.
\end{proof}
Consequently, constructing a label/default adapter and asserting a contract are
different operations. Keeping that distinction also preserves associativity.
A full experimental compatibility layer can implement its original invocation
relation as a checked adapter, while translating its public guarantees into the
appropriate quantified predicates. Strict kind and callable checks are verified
at that boundary rather than silently deferred by the translation.

\section{Modular guarantees}
Universal constraints express a statement the call-restriction interpretation
does not express by itself: one selected implementation preserves its result type at every
admitted instance whenever that instance returns successfully. This supports checking \code{map}, polymorphic callbacks, and bounded
refinement-preserving utilities without seeing their eventual calls. The
intersection algebra preserves input/output correlations. Dependent universals
state relationships across argument and result values. Existential seals support
representation-independent replacement under an explicit theorem.

Untyped callbacks and struct encodings can express many of the same
computations. Quantified descriptions additionally preserve their relationships
as constraints.

\section{CDuce and semantic subtyping}
Frisch, Castagna, and Benzaken develop semantic subtyping with function and
Boolean types; Castagna and Xu extend its foundation to polymorphic types and
convex models \cite{semantic,castagnaxu}. CDuce's polymorphic calculi combine
instantiation with intersections and preserve input/output correlations
\cite{poly1,poly2}. These supply close precedents for the arrow and instantiation
layers here.

The current proposal adds explicit quantifiers to the general CUE predicate
language. Its partial-correctness arrows refer to a fixed operational
inhabitant model, with impredicative quantification over sets of inhabitants. A subtyping algorithm proved for a different arrow model must
be revalidated for the selected reference fragment. Convexity of a semantic
subtyping model and relational parametricity are separate properties.

\section{Elixir and checked dynamic execution}
Elixir's set-theoretic work uses unions, intersections, refinements, and guarded
function behavior. The strong-arrow account pays particular attention to
behavior outside the input domain and interaction with dynamic values
\cite{elixirstrong}. It is relevant to CUE's bridge and legacy-expression
boundaries: failure behavior deserves an explicit interpretation alongside the
successful result type.

This proposal's universals are general constraint binders. Its partial-correctness and optional effect
arrows are specified by \cref{ch:functions}, rather than identified with Elixir's
strong-arrow construction. Both approaches use set-theoretic descriptions to
retain more precise information than an uncorrelated union of inputs and outputs.

\section{Hyperdoctrines and compact closure}
Consider objects that are finite interfaces and morphisms that are definable
relations $R\subseteq X\times Y$. Composition is
\[
 (S\circ R)(x,z)\equiv\exists y.\;R(x,y)\land S(y,z).
\]
Juxtaposition of interfaces is the tensor. Equality gives the unit/counit
relations making every object self-dual; the snake identities follow by
existentially eliminating equality witnesses. This is the compact structure of
a relational category.

Adding universal formulas enriches the class of definable relations while
retaining these operations, provided the fragment is closed under the required
composition and equality relations. Conjunction within one predicate fiber is
distinct from the tensor of interfaces. Executable functions form a
separate operational subcollection. The compact structure belongs to the
relations and does not require executable total functions.

Lawvere's adjunctions and the bifibrational account of Melli\`es and Zeilberger
organize these refinement predicates over the relational interface category
\cite{lawvere,mellies}.

\chapter{Requirements and implementation outcome}
\label{ch:requirements}
The addendum's section 9 lists six requirements. The last is its restriction on
expressive strength, referring to section 7. The first five are evaluated below
against the defined reference semantics \cite{cueaddendum}.

\section{Directionality}
A reusable capability is a universal predicate on the implementation. A
particular call is a relation on existential argument and result cells. Neither
judgment guesses whether a field is provided or consumed. A higher-order
parameter introduces its required capability in the caller's type. Module
replacement uses its explicitly declared interface and observation relation.

\section{Monotone refinement}
Conjunction refines the joint predicate. Both quantifiers preserve that
refinement, by \cref{thm:monotone}. Dependency-preserving copying and retention
of bound variables prevent accidental weakening of universal obligations.
An established contradiction remains empty, and an atomic concrete result can
only persist or conflict. \Cref{ch:residual} implements this property by retaining
mandatory annotation certificates, explicit source assertions, and
validation dependencies. Relevant-consistency errors block explicit interfaces;
they do not rewrite inhabited function predicates to bottom. Annotation erasure preserves the source body, and
its declaration predicates obey \cref{eq:declaration-homomorphism}. Defaults and comprehension execution retain the
separate observation treatment already identified in the reference model.

\section{Cross-file input and output invariants}
Existing argument/result records express existential call restrictions.
Universally quantified declarations express reusable behavior and dependent
relations. Both can be conjoined across files. \cref{ch:programs,ch:dependent}
give concrete examples of each form, including length preservation and
request-indexed responses.

\section{Callable self-unification}
Function inhabitants have stable semantic origins and captured environments.
Their equality is structural on these descriptors, rather than pointer-based.
Conjunction is idempotent on their singleton predicates, including closures
reached through copying or unrelated record refinement. Partial applications
use normalized original-slot bindings. Sealed operation handles preserve the
interface's public alias structure.

\section{Idiomatic foreign interfaces}
Portable logical domains remain in public signatures. Representability and
conversion failures may become ordinary bottom; an explicit effect records them
when their tags are observable. The existing bridge performs its existing
checks and reports those outcomes. Logical kind/domain mismatches
are distinguished from in-domain foreign failures. Native widths remain private
bridge metadata. No additional conversion stage is required by quantification.

\section{Expressive strength and implementation milestones}
The reference semantics admits higher-rank and dependent specifications. An
execution profile may retain the current nonrecursive policy or admit
recursion with partial-correctness typing. Neither checking nor arrow membership
requires a termination proof. \Uone{} is an implementation entry point with explicit generic
capabilities and semantic arrow intersections. \Sone{} adds abstract type witnesses
at module boundaries. \Shigher{} extends the same formation rules to arbitrary
well-sorted positions.

A conforming implementation establishes the following contracts:
\begin{enumerate}
\item Elaboration preserves legacy data observations, lexical binding, copied
witness sharing, and quantifier dependencies.
\item Every explicit interface passes relevant consistency, independently of
implementation linkage. The checker verifies every mandatory strict judgment
before accepting executable code. Only explicit source operations and data hypotheses may
remain pending; unproved annotation implications block certification. Proved
empty predicates normalize to bottom at their retained scopes.
\item Symbolic evaluation preserves the successful-completion semantics and the
specified observation/effect behavior.
\item The abstraction boundary preserves seal identity, public aliasing, and
the equality-respecting simulation used for module evolution.
\end{enumerate}

\section{Conclusion}
Universal and existential binders are dual projections of the same relation-valued
constraint semantics. Unification remains intersection. Function types arise by
universally constraining an operational application relation; their variance and
polymorphic behavior follow from that definition. Dependent contexts and
existential modules extend the same structure. The resulting language can be
implemented in rank-restricted stages without changing what its quantified
descriptions mean. Relevant consistency checks the successful observations
claimed by explicit interfaces, while partial-correctness typing preserves
ordinary computation failure.

\appendix
\chapter{Core judgments and finite regression model}
\label{app:core}
\section{Formation and binding summary}
A reference type expression has the grammar
\[
\begin{aligned}
 T,U ::= {}& b\mid a\mid T\land U\mid T\lor U\mid\neg T
       \mid T\times U\mid\List(T)\mid T\Rightarrow U\\
       &\mid \{\ell_i:T_i\}_{i\in I}
       \mid \forall a:S\,T\mid\exists a:S\,T.
\end{aligned}
\]
Here $b$ is a base predicate; $a$ is a context variable of the appropriate sort;
$S$ is a well-formed value sort or the type sort $\Type$, possibly with a guard.
Complement is relative to the subject sort. A surface profile may retain only
its supported decidable negative predicates. Dependent descriptions extend $T_i$
with references to earlier argument or record cells, using the shared binding
graph. Function effects refine the application predicate of \cref{ch:functions}.

Contexts are telescopes. The free variables of a binder's sort and guard must
already occur in its context. A bound variable is identified by its binder, not
by its spelling. A quantified expression keeps its subject outside the new
binder; an explicit dependent sum keeps the witness as part of the subject.

\section{Sound introduction and elimination principles}
Write $\Gamma;\Delta\vdash e:T$ for a typing judgment in a term environment
$\Gamma$ and a rigid constraint context $\Delta$. The type variable $a$ is fresh
for $\Gamma$ in universal introduction:
\[
\frac{\Gamma;\Delta,a:S\vdash e:T}
     {\Gamma;\Delta\vdash e:\forall a:S\,T}
\qquad
\frac{\Gamma;\Delta\vdash e:\forall a:S\,T\qquad\Delta\vdash s:S}
     {\Gamma;\Delta\vdash e:T[s/a]}.
\]
The first premise checks the same executable expression at an arbitrary
admissible assignment. Bounds appear as assumptions in $\Delta,a:S$.

For transparent existential descriptions, introduction chooses a witness and
elimination checks the continuation at an arbitrary opened witness:
\[
\frac{\Gamma;\Delta\vdash e:T[s/a]\qquad\Delta\vdash s:S}
     {\Gamma;\Delta\vdash e:\exists a:S\,T}.
\]
Elimination preserves the subject and shares the witness across its projections.
Its result may not expose the local witness without an enclosing existential.
Opaque packages additionally use the sealing discipline in \cref{ch:modules}.

Intersection introduction uses the same term:
\[
\frac{\Gamma;\Delta\vdash e:T\qquad\Gamma;\Delta\vdash e:U}
     {\Gamma;\Delta\vdash e:T\land U}.
\]
Source-written interfaces additionally pass relevant consistency at their
retained roots. Inferred types and instantiated proof conclusions remain in the
full semantic type algebra, including empty successful-result predicates.
Subsumption uses a checked strict inclusion derivation. Application must pass
its static callable and argument judgments. Its conclusion concerns successful
results, not termination or absence of bottom. A runtime refinement operation
must occur in the source, with its rule checked as in
\cref{sec:dynamic-admission}. All annotation implications are proved before
certification; their proofs erase.

For returned values, the introduction rules give membership in the
corresponding intersection or union; computation judgments quantify only over
successful returns. Elimination uses the selected instance or a scoped
existential witness. These principles specify the reference calculus independently
of a particular constraint-solving strategy.

\section{Regression checks}
The source archive includes three executable standard-library Python models.

\code{checks/check_model.py} enumerates a finite partial-function algebra and
checks arrow laws, partial identities, failure-only call regions,
quantifier laws, scope conditions, and set-valued meet. A partial identity witnesses the nonempty denotation of the rejected
\code{quiet} interface. The model also checks a finite
applicative model whose type domain is the full powerset of its value carrier;
quantified types can occur as instance arguments in that model.

\code{checks/check_kernel.py} implements a small certificate-oriented structural
checker and the static/dynamic gate. Its fixtures cover singleton and kind
inclusions, record width/depth and closed field lookup, union/intersection
rules, explicit impredicative substitution, the generic meet body, blocked
annotation obligations, and explicit dependent assertions. Declaration
fixtures check regrouping and common erasure for the five presentations.
Negative fixtures verify that unproved result bounds, wrong-kind results, and
exhausted static budgets produce no executable filters.

\code{checks/check_relevance.py} checks explicit-domain observation partitions,
incompatible overlaps, multiway result conflicts, conjunction coherence,
source-anchored universal instances, and the distinction between generic
interface admission and derived failure types. Its fixtures also check that
relevance rejection is independent of denotational emptiness and produces no
executable filter. The exact counts are recorded in \code{checks/results.json}.

The models test finite instances and selected kernel rules. The implementation
contract is the rule set and conformance table in \cref{ch:residual}; complete
integration with CUE's graph evaluator requires the elaboration and demand
invariants in \cref{ch:checking}. In particular, finite tests are not a proof of
that integration or of complete semantic-subtyping inference.

\chapter{Example and feature index}
\label{app:index}
\begin{longtable}{@{}P{0.35\textwidth}P{0.59\textwidth}@{}}
\toprule
Feature & Location and program forms\\\midrule
\endfirsthead
\toprule Feature & Location and program forms\\\midrule
\endhead
Binding syntax & \Cref{ch:surface}: quantified expressions, declaration
parameters, block prefixes, function-local generics, parametric aliases, and
instantiation selection.\\
Pointwise unification & \Cref{ch:refinement}: alpha-renaming, shared subject,
constant constraints, guarded bounded clauses, non-distributivity.\\
Rank-1 programming & \Cref{ch:programs}: identity, bounded pairing, minimum,
map, filter, partition, flat mapping, composition, adapters.\\
Higher-rank programming & \Cref{ch:programs}: a polymorphic callback used at two
types, returned polymorphic closures, quantified utility records.\\
Calling conventions & \Cref{ch:programs}: all parameter modes, defaults,
omission, hidden labels, partial application, open protocols, attributes.\\
Dependent contracts & \Cref{ch:dependent}: exact successor, interval clamp,
length-preserving map, safe indexing, append, zip, split, replication.\\
Dependent structures & \Cref{ch:dependent}: dimension-carrying matrices,
transpose, dot product, matrix multiplication, indexed service responses.\\
Existential reuse & \Cref{ch:modules}: sealed counters, generic clients, module
laws, heterogeneous packages, shared-witness comparison, codecs.\\
Compatibility & \Cref{ch:comparison}: feature-by-feature correspondence,
adapters, old data code, result combination, experimental migration.\\
Semantic guarantees & \Cref{ch:quantifiers,ch:refinement,ch:functions}:
adjunctions, floating, monotonicity, partial correctness, variance and application.\\
Mandatory checking & \Cref{ch:residual}: strict System F/structural certificates,
scoped eager refutation, relevant consistency of explicit observations,
rigid declaration checks, explicit assertions, declaration coherence,
annotation erasure, stage gates, and resource bounds.\\
Evolution guarantees & \Cref{ch:modules}: equality-respecting representation
simulation, quantified observations, old-compatible views.\\\bottomrule
\end{longtable}

\clearpage
\renewcommand{\bibname}{References}
\begin{thebibliography}{99}
\small
\addcontentsline{toc}{chapter}{References}
\raggedright

\bibitem{cuespec}
The CUE Authors.
\emph{CUE language specification}.
Living specification, accessed 21 September 2026.
\url{https://cuelang.org/docs/reference/spec/}.

\bibitem{cuerefs}
The CUE Authors.
\emph{CUE language specification: References}.
Specification of copying and rebinding referenced expressions.
Accessed 21 September 2026.
\url{https://cuelang.org/docs/reference/spec/#references}.

\bibitem{cuedefaults}
The CUE Authors.
\emph{Specifying a default value for a field}.
Accessed 21 September 2026.
\url{https://cuelang.org/docs/howto/specify-a-default-value-for-a-field/}.

\bibitem{cuelist}
The CUE Authors.
\emph{Package list}.
Standard-library interface documentation, accessed 21 September 2026.
\url{https://pkg.go.dev/cuelang.org/go/pkg/list}.

\bibitem{cueproposal}
Marcel van Lohuizen.
\emph{Proposal: functions in CUE}. 10 July 2026.
CUE proposal 4484, draft at repository revision
\texttt{a949162aedb86661183c46abef47919d16605014}.
\url{https://github.com/cue-lang/proposal/blob/a949162aedb86661183c46abef47919d16605014/designs/language/4484-functions.md}.

\bibitem{cueaddendum}
Marcel van Lohuizen.
\emph{Why CUE Functions Tighten Instead of Widen}.
Theory addendum to proposal 4484, committed 19 September 2026,
repository revision \texttt{a949162aedb86661183c46abef47919d16605014}.
\url{https://github.com/cue-lang/proposal/blob/a949162aedb86661183c46abef47919d16605014/designs/language/4484-theory-addendum.md}.

\bibitem{lawvere}
F. William Lawvere.
Adjointness in foundations.
\emph{Dialectica}, 23(3--4):281--296, 1969.
Reprinted with commentary in \emph{Reprints in Theory and Applications of
Categories}, 16:1--16, 2006.
\url{https://tac.mta.ca/tac/reprints/articles/16/tr16abs.html}.

\bibitem{mellies}
Paul-Andr\'e Melli\`es and Noam Zeilberger.
\emph{A bifibrational reconstruction of Lawvere's presheaf hyperdoctrine}.
2016. arXiv:1601.06098.
\url{https://arxiv.org/abs/1601.06098}.

\bibitem{reynolds84}
John C. Reynolds.
Polymorphism is not set-theoretic.
In \emph{Semantics of Data Types}, LNCS 173, pages 145--156, 1984.
\url{https://doi.org/10.1007/3-540-13346-1_7}.

\bibitem{reynolds83}
John C. Reynolds.
Types, abstraction and parametric polymorphism.
In \emph{Information Processing 83}, pages 513--523. North-Holland, 1983.
\url{https://www.cs.cmu.edu/~crary/819-f09/Reynolds83.pdf}.

\bibitem{semantic}
Alain Frisch, Giuseppe Castagna, and V\'eronique Benzaken.
Semantic subtyping: Dealing set-theoretically with function, union, intersection,
and negation types.
\emph{Journal of the ACM}, 55(4), Article 19, 2008.
\url{https://doi.org/10.1145/1391289.1391293}.

\bibitem{castagnaxu}
Giuseppe Castagna and Zhiwu Xu.
Set-theoretic foundation of parametric polymorphism and subtyping.
In \emph{ICFP}, pages 94--106, 2011.
\url{https://doi.org/10.1145/2034773.2034788}.

\bibitem{poly1}
Giuseppe Castagna, Kim Nguyen, Zhiwu Xu, Hyeonseung Im, Sergue\"i Lenglet,
and Luca Padovani.
Polymorphic functions with set-theoretic types. Part 1: Syntax, semantics,
and evaluation.
In \emph{POPL}, pages 5--18, 2014.
\url{https://doi.org/10.1145/2535838.2535840}.

\bibitem{poly2}
Giuseppe Castagna, Kim Nguyen, Zhiwu Xu, and Pietro Abate.
Polymorphic functions with set-theoretic types. Part 2: Local type inference
and type reconstruction.
In \emph{POPL}, pages 289--302, 2015.
\url{https://doi.org/10.1145/2676726.2676991}.

\bibitem{cducetutorial}
The CDuce Project.
\emph{Polymorphism tutorial}.
Accessed 21 September 2026.
\url{https://www.cduce.org/tutorial_polymorphism.html}.

\bibitem{cln2024}
Giuseppe Castagna, Micka\"el Laurent, and Kim Nguyen.
Polymorphic type inference for dynamic languages.
\emph{Proceedings of the ACM on Programming Languages}, 8(POPL), 2024.
\url{https://doi.org/10.1145/3632882}.
Preprint: \url{https://arxiv.org/abs/2311.10426}.

\bibitem{elixirstrong}
Jos\'e Valim.
\emph{Strong arrows: a new approach to gradual typing}.
20 September 2023.
\url{https://elixir-lang.org/blog/2023/09/20/strong-arrows-gradual-typing/}.

\bibitem{mitchellplotkin}
John C. Mitchell and Gordon D. Plotkin.
Abstract types have existential type.
\emph{ACM Transactions on Programming Languages and Systems},
10(3):470--502, 1988.
\url{https://doi.org/10.1145/44501.45065}.

\bibitem{mitchell86}
John C. Mitchell.
Representation independence and data abstraction.
In \emph{POPL}, pages 263--276, 1986.
\url{https://doi.org/10.1145/512644.512669}.

\bibitem{dunfield}
Jana Dunfield and Neelakantan R. Krishnaswami.
Complete and easy bidirectional typechecking for higher-rank polymorphism.
In \emph{ICFP}, pages 429--442, 2013.
\url{https://doi.org/10.1145/2500365.2500582}.
Preprint: \url{https://arxiv.org/abs/1306.6032}.

\bibitem{nadathur}
Gopalan Nadathur, Bharat Jayaraman, and Keehang Kwon.
\emph{Scoping constructs in logic programming: Implementation problems
and their solution}.
1998. arXiv:cs/9809016.
\url{https://arxiv.org/abs/cs/9809016}.

\bibitem{liquid}
Patrick M. Rondon, Ming Kawaguchi, and Ranjit Jhala.
Liquid types.
In \emph{PLDI}, 2008.
\url{https://goto.ucsd.edu/~rjhala/papers/liquid_types.html}.

\bibitem{ghcrank}
The GHC Developers.
\emph{GHC User's Guide: Arbitrary-rank polymorphism}.
Accessed 21 September 2026.
\url{https://ghc.gitlab.haskell.org/ghc/doc/users_guide/exts/rank_polymorphism.html}.


\bibitem{cuecycles}
The CUE Authors.
\emph{Reference Cycles}.
CUE language tour, accessed 22 September 2026.
\url{https://cuelang.org/docs/tour/references/cycle/}.

\bibitem{cousot77}
Patrick Cousot and Radhia Cousot.
Abstract interpretation: A unified lattice model for static analysis of programs
by construction or approximation of fixpoints.
In \emph{POPL}, pages 238--252, 1977.
\url{https://www.di.ens.fr/~cousot/COUSOTpapers/POPL77.shtml}.

\bibitem{ccp91}
Vijay A. Saraswat, Martin Rinard, and Prakash Panangaden.
The semantic foundations of concurrent constraint programming.
In \emph{POPL}, pages 333--352, 1991.
\url{https://doi.org/10.1145/99583.99627}.


\bibitem{pistone}
Paolo Pistone.
On completeness and parametricity in the realizability semantics of System F.
\emph{Logical Methods in Computer Science}, 15(4), Article 6, 2019.
\url{https://doi.org/10.23638/LMCS-15(4:6)2019}.
Preprint: \url{https://arxiv.org/abs/1802.05143}.

\bibitem{bidirsurvey}
Jana Dunfield and Neel Krishnaswami.
Bidirectional typing.
\emph{ACM Computing Surveys}, 54(5), 2021.
\url{https://doi.org/10.1145/3450952}.
Preprint: \url{https://arxiv.org/abs/1908.05839}.

\bibitem{jiang2025}
Shengyi Jiang, Chen Cui, and Bruno C. d. S. Oliveira.
Bidirectional higher-rank polymorphism with intersection and union types.
\emph{Proceedings of the ACM on Programming Languages}, 9(POPL),
Article 71, pages 2118--2148, 2025.
\url{https://doi.org/10.1145/3704907}.

\bibitem{explicitrefine}
Jad Elkhaleq Ghalayini and Neel Krishnaswami.
\emph{Explicit refinement types}.
2023. arXiv:2311.13995.
\url{https://arxiv.org/abs/2311.13995}.

\bibitem{nonstrict}
Tommaso Petrucciani, Giuseppe Castagna, Davide Ancona, and Elena Zucca.
Semantic subtyping for non-strict languages.
In \emph{TYPES 2018}, LIPIcs 130, Article 4, 2019.
\url{https://doi.org/10.4230/LIPIcs.TYPES.2018.4}.
Preprint: \url{https://arxiv.org/abs/1810.05555}.

\end{thebibliography}
\end{document}
